\chapter{Model forms and semipositivity}
\index{model form}
\index{semipositivity!BFJ}
\index{relative nefness}
\label{chap:bfj-model-forms}

The differential-geometric objects of the BFJ theory are assembled from
algebraic data on models.  A closed form is a numerical class on some model,
a potential is a vertical divisor, and positivity means relative nefness.
The purpose of this chapter is to make that dictionary precise and to record
the intersection calculations that control all later compactness and energy
arguments.

\section{Orientation}
\label{sec:ch20-orientation}

\subsection{What has to be built, and why}
\label{sec:ch20-why}

\Cref{chap:metrics} produced a supply of concrete objects: model metrics on a
line bundle, their relative weights, the model functions
\(\varphi_D\) attached to vertical divisors, and the Chambert--Loir measures
obtained from special-fibre intersection numbers.  All of these are attached
to a \emph{choice} of model.  A metric that is defined on \(\cX\) is also
defined on every model dominating \(\cX\), and the two descriptions carry the
same analytic information; but they are not literally the same algebraic
datum, and no single model carries all of them.

The pluripotential theory that follows cannot be organized model by model.
Its central operations---taking a supremum of a family of potentials,
regularizing from above by a decreasing net, solving a Monge--Amp\`ere
equation---leave every fixed model immediately.  What is needed is a
\emph{category of curvature data} that is closed under these operations at
the linear level, in which one can add, subtract, and compare without asking
first on which model the answer lives.

The construction that supplies it is elementary and is the subject of
\Cref{sec:ch20-numerical,sec:ch20-model-forms}.  For each model \(\cX\) one
forms the finite-dimensional real vector space \(N^1(\cX/R)\) of line-bundle
classes modulo numerical equivalence on vertical curves; a domination
\(\rho:\cY\to\cX\) induces an injection \(\rho^*\), and the direct limit
\[
 \mathcal Z^{1,1}(X)=\varinjlim_{\cX}N^1(\cX/R)
\]
is the space of \emph{closed model \((1,1)\)-forms}.  Three separate things
then have to be proved, and they organize the chapter.

\begin{enumerate}
\item\label{itm:ch20-plan:limit}
 \emph{The limit has to be shown to compute something.}  A direct limit over
 an index category is only as good as its transition maps.  We prove that
 the models form a directed set (\Cref{prop:common-domination}), that
 \(\rho^*\) is injective (\Cref{lem:ch20-pullback-injective}), and that two
 determinations define the same form precisely when they agree on
 \emph{one} common domination
 (\Cref{prop:ch20-two-model-comparison}).  Without the last statement,
 ``the same form on two models'' would be an assertion about infinitely many
 models rather than a finite computation.
\item\label{itm:ch20-plan:ddc}
 \emph{Potentials have to be attached to forms.}  A vertical divisor \(D\)
 carries both a function \(\varphi_D\) on \(X^{\An}\) and a numerical class
 \([D]\); the operator \(\ddc\) of \Cref{sec:ch20-ddc} is the passage from
 the first to the second.  Its kernel, image, and cokernel are the content
 of the exact sequence \Cref{thm:ch20-ddc-exact}, which is the single
 structural theorem of the chapter and the reason all later normalizations
 close up.
\item\label{itm:ch20-plan:positivity}
 \emph{Positivity has to be defined, and the competing definitions
 separated.}  There are four candidate notions---relative ampleness,
 relative nefness, numerical positivity on the special fibre, and nefness of
 the generic class---and they are genuinely different.
 \Cref{sec:ch20-positivity} proves every implication that holds and gives an
 explicit counterexample for every implication that fails.
\end{enumerate}

Finally, everything downstream is quantitative, so the chapter ends with the
two intersection-theoretic inputs that the analysis will use without further
comment: the intersection pairing with one vertical-divisor factor
(\Cref{sec:ch20-intersection}) and the relative Hodge index inequality
(\Cref{thm:ch20-mixed-hodge}) that makes energy differences concave.

\subsection{The payoff, stated in advance}
\label{sec:ch20-payoff}

Four statements proved below are used continually in
\Crefrange{chap:bfj-psh}{chap:bfj-examples}, and the reader may wish to see
them before the constructions.

\begin{enumerate}
\item \emph{Every closed model form with ample generic class becomes
 relatively ample, on a cofinal family of models, after adding a model
 potential} (\Cref{prop:ch20-positive-representative}).  When the form is
 already semipositive, the correcting potential may be taken arbitrarily
 small.  This is what allows the whole theory to be developed with a
 \emph{positive} reference form, which in turn is what makes constants
 \(\theta\)-psh and gives the maximum construction a base point.
\item \emph{The model differential has kernel exactly the constants and
 image exactly the forms with numerically trivial generic class}
 (\Cref{thm:ch20-ddc-exact}).  In particular, a closed model form is
 determined by its generic numerical class together with a model function,
 up to an additive constant.  Every later reduction of the form ``replace
 \(\theta\) by \(\theta+\ddc\psi\) and \(\varphi\) by \(\varphi-\psi\)''
 rests on this.
\item \emph{The \(\theta\)-psh model functions are stable under finite
 maxima} (\Cref{prop:ch20-model-max}).  This is not a formal consequence of
 convexity of the nef cone; the proof passes through relative global
 generation of a perturbed ideal, and the reason a naive argument fails is
 isolated in \Cref{rem:ch20-why-max-is-hard} and made explicit in
 \Cref{ex:ch20-max-fails-naively}.
\item \emph{Vertical divisors have nonpositive self-intersection against
 relatively nef classes, with a sharp equality case}
 (\Cref{prop:ch20-relative-hodge}, \Cref{thm:ch20-mixed-hodge}).  This is
 the algebraic origin of concavity of the Monge--Amp\`ere energy in
 \Cref{chap:bfj-energy} and of the Cauchy--Schwarz estimates in
 \Cref{chap:bfj-capacity,chap:bfj-finite-energy-measures}.
\end{enumerate}

A fifth statement is negative and is worth flagging with the others.  A
vertical class is relatively nef only if it is a multiple of the special
fibre (\Cref{thm:ch20-vertical-nef-rigidity}); equivalently, the only
\(0\)-psh model functions are the constants.  Everything interesting in the
theory therefore happens because \(\theta\) is present, not because the
vertical divisors themselves carry positivity.

\subsection{Standing hypotheses: the original BFJ signature}
\label{sec:ch20-hypotheses}

Throughout \Crefrange{chap:bfj-model-forms}{chap:bfj-continuity}, unless an
extension is explicitly labelled, the following hypotheses are in force.  We
refer to them as the \emph{original BFJ signature} and abbreviate them
\hyperref[eq:ch20-standing]{\textup{(BFJ)}}.

\begin{equation}
\label{eq:ch20-standing}
\left\{
\begin{aligned}
 &\text{\(K\) is complete, \emph{discretely} valued, with residue field
   \(k\) of characteristic \(0\);}\\
 &\text{\(R=K^\circ\), \(\pi\in R\) a uniformizer, and \(v_K(\pi)=1\);}\\
 &\text{\(X\) is a smooth, connected, projective \(K\)-variety of dimension
   \(n\geq1\);}\\
 &\text{models \(\cX\) are normal, flat, and projective over \(R\) with
   \(\cX_K\simeq X\);}\\
 &\text{domination is a projective birational \(R\)-morphism inducing the
   identity on \(X\).}
\end{aligned}
\right.
\end{equation}

These are the hypotheses of \cite{BFJ16}, and we keep them visible rather
than absorbing them into the notation.  \Cref{sec:ch20-scope} records, statement by statement, which of
them each result actually uses and which are known to be removable.

Two normalizations are used constantly and are quoted here so that no proof
below has to rederive them.  The first is the sign of a model function
(\Cref{prop:metrics-divisor-model-function}, \eqref{eq:appJ-model-divisor}):
if \(D\) is a vertical Cartier divisor with local equation \(f_D\), then
\(\cO(D)\) is locally generated by \(f_D^{-1}\) and

\begin{equation}
\label{eq:ch20-sign}
              \varphi_D(x)=-\log|f_D(x)|=v_x(f_D),
              \qquad \varphi_{\mathcal X_s}=1.
\end{equation}

Thus effective vertical divisors give nonnegative model functions, and the
special fibre gives the constant \(+1\).  The second is the
Chambert--Loir pairing of \Cref{prop:metrics-cl-model-function-integral},
\eqref{eq:metrics-cl-pairing}, which converts a vertical intersection number
into an integral of a model function.  Together with the projection formula
\Cref{thm:model-projection-formula} these two facts generate every sign in
the chapter, and \Cref{sec:ch20-normalization-audit} checks them against the
unit tests of \Cref{sec:appJ-unit-tests}.

\subsection{What is proved here and what is imported}
\label{sec:ch20-imported}

Everything in \Crefrange{sec:ch20-numerical}{sec:ch20-intersection} is
proved, with the following exceptions.  Each is quoted with its hypotheses
at the point of use, and each is a statement of classical algebraic
geometry rather than of pluripotential theory.

\begin{enumerate}
\item\label{itm:ch20-import:res}
 \emph{Resolution of singularities in characteristic zero}, in the form
 already used in \Cref{chap:skeleta}: SNC models are cofinal among models
 (\Cref{prop:snc-cofinal}).  This is where the hypothesis
 \(\operatorname{char}k=0\) enters, and it enters only here and in the
 statements that quote it.
\item\label{itm:ch20-import:ega}
 \emph{Openness of relative ampleness in the fibres},
 \cite[4.7.1]{EGAIII1}, used once, in
 \Cref{lem:ch20-special-fibre-test}, to reduce relative positivity over
 \(R\) to ordinary positivity on the projective \(k\)-scheme \(\cX_s\).
\item\label{itm:ch20-import:kleiman}
 \emph{Kleiman's theorem} that the nef cone of a projective scheme over a
 field is the closure of the ample cone, and that products of nef classes
 have nonnegative degree on effective cycles; and the
 \emph{Nakai--Moishezon criterion} for a \(\mathbf Q\)-class.  Both are used
 only through \Cref{lem:ch20-special-fibre-test}, that is, on the projective
 \(k\)-scheme \(\cX_s\), where they are the standard statements.
\item\label{itm:ch20-import:kunnemann}
 \emph{K\"unnemann's extension argument} for numerically trivial line
 bundles, together with the Chow-theoretic projection formula.  This is the
 one genuinely nonformal input to \Cref{thm:ch20-ddc-exact}, and it is
 isolated as \Cref{lem:ch20-kunnemann-input} so that the reader can see
 exactly how much is being assumed.
\item\label{itm:ch20-import:extension}
 \emph{Existence of a model on which a given ample generic class extends to
 a relatively ample class}, which for a projective \(X\) is
 \Cref{prop:line-bundle-extension} together with the flat-closure
 construction of \Cref{chap:models}.
\item\label{itm:ch20-import:blowup}
 \emph{Domination of a modification by a blowup},
 \Cref{thm:appB-admissible-domination}, and \emph{relative ampleness of a
 twist by a large multiple}, \cite[4.6.13]{EGAII}.  These two are used only
 in \Cref{lem:ch20-ample-propagation}, which is the mechanism by which
 relative ampleness with a prescribed generic class is transported upward
 through the inverse system of models.
\end{enumerate}

No analytic input is used in this chapter.  In particular nothing from
\Cref{part:cld} is needed: the comparison with the Chambert--Loir--Ducros
calculus is a theorem of \Cref{chap:cld-examples}
(\Cref{thm:ch19-cld-bfj-comparison}), not a definition, and we are careful in
\Cref{sec:ch20-normalization-audit} to keep the two operators \(\ddc\) and
\(\dprimeddprime\) typographically and conceptually apart.

\subsection{Plan}
\label{sec:ch20-plan}

\Cref{sec:ch20-numerical} constructs the relative numerical spaces
\(N^1(\cX/R)\), proves that they are finite dimensional, that pullback along
a domination is injective, and that all positivity questions may be tested
on the special fibre.  \Cref{sec:ch20-model-forms} performs the direct
limit, proves the two-model comparison criterion, and shows that
semipositivity is independent of the determination.
\Cref{sec:ch20-ddc} constructs the model differential, verifies its
normalization against \Cref{app:normalizations}, and computes it explicitly
on a curve.  \Cref{sec:ch20-exact} proves the exact sequence and its
corollaries, and \Cref{sec:ch20-positive-representative} deduces from it
that a form with ample generic class has an \(\mathcal X\)-positive
representative.  \Cref{sec:ch20-positivity} separates the four positivity
notions.  \Cref{sec:ch20-max-envelopes} treats vertical ideals, the maximum
construction, and model envelopes.  \Cref{sec:ch20-intersection} constructs
the intersection pairings and proves the Hodge-type inequalities.
\Cref{sec:ch20-examples} carries five examples to explicit numbers,
\Cref{sec:ch20-scope} is the scope table separating the original BFJ theory
from its later extensions, and \Cref{sec:ch20-dictionary} is the chapter
dictionary.

\section{Relative numerical classes on a model}
\label{sec:ch20-numerical}

\subsection{The space \texorpdfstring{\(N^1(\cX/R)\)}{N1(X/R)}}
\label{sec:ch20-n1}

Fix a model \(\cX\) as in \eqref{eq:ch20-standing}.  Its special fibre
\(\cX_s\) is a projective scheme over \(k\), of pure dimension \(n\) because
\(\cX\to\Spec R\) is flat and projective of relative dimension \(n\).  A
\emph{vertical curve} on \(\cX\) is an integral closed subscheme
\(C\subseteq\cX_s\) of dimension one; equivalently, it is an integral curve
in \(\cX\) whose image under \(\cX\to\Spec R\) is the closed point.  This
condition excludes horizontal curves, which dominate \(\Spec R\).

\begin{definition}[Relative numerical equivalence]
\index{numerical class!relative}\index{vertical curve}
\label{def:ch20-relative-numerical}
Two classes \(\alpha,\beta\in\Pic(\cX)_{\mathbf R}\) are
\emph{relatively numerically equivalent}, written
\(\alpha\equiv_{\cX_s}\beta\), if \((\alpha\cdot C)=(\beta\cdot C)\) for
every vertical curve \(C\).  The quotient is written

\begin{equation}
\label{eq:ch20-n1-definition}
 N^1(\cX/R):=\Pic(\cX)_{\mathbf R}/\equiv_{\cX_s},
\end{equation}

and the class of \(\cL\) in it is denoted \(c_1(\cL)\) or
\([\cL]\).  We call \(N^1(\cX/R)\) the space of \emph{relative numerical
classes}.
\end{definition}

The definition uses only degrees on curves contracted by
\(\cX\to\Spec R\), namely the vertical curves.  Horizontal curves may be
proper over \(R\), but they dominate the base and do not enter relative
numerical equivalence.  This is the precise sense in which
\(N^1(\cX/R)\) is a \emph{relative} invariant.

\begin{lemma}
\label{lem:ch20-special-fibre-test}
Let \(\cX\) be a model and \(\cL\) a line bundle on \(\cX\), with
restriction \(\cL_s:=\cL|_{\cX_s}\).
\begin{enumerate}
\item\label{itm:ch20-sft:inject} The restriction map
 \(\Pic(\cX)_{\mathbf R}\to N^1(\cX_s)\) to the numerical class group of the
 projective \(k\)-scheme \(\cX_s\) induces an injection
 \(N^1(\cX/R)\hookrightarrow N^1(\cX_s)\).  In particular
 \(N^1(\cX/R)\) is a finite-dimensional real vector space.
\item\label{itm:ch20-sft:nef} \(\cL\) is relatively nef, that is
 \((\cL\cdot C)\geq0\) for every vertical curve \(C\), if and only if
 \(\cL_s\) is nef on \(\cX_s\).
\item\label{itm:ch20-sft:ample} \(\cL\) is relatively ample over \(R\) if
 and only if \(\cL_s\) is ample on \(\cX_s\).
\item\label{itm:ch20-sft:cone} The relatively nef classes form a closed
 convex cone in \(N^1(\cX/R)\) with nonempty interior, and it is the closure
 of the open convex cone of relatively ample classes.  A sum of a relatively
 nef and a relatively ample class is relatively ample.
\end{enumerate}
\end{lemma}

\begin{proof}
\emph{\cref{itm:ch20-sft:inject}.}  By construction the kernel of
\(\Pic(\cX)_{\mathbf R}\to N^1(\cX_s)\) consists of the classes with zero
degree on every curve in \(\cX_s\), which is exactly \(\equiv_{\cX_s}\).
Finite dimensionality of \(N^1(\cX_s)\) for a projective scheme over a field
is the theorem of the base in the form recalled in
\Cref{def:appF-numerical}.

\emph{\cref{itm:ch20-sft:nef}.}  Immediate from the definition: vertical
curves on \(\cX\) are exactly the curves on \(\cX_s\), and degrees are
computed on the closed subscheme.

\emph{\cref{itm:ch20-sft:ample}.}  If \(\cL\) is relatively ample then its
restriction to the fibre \(\cX_s\) is ample.  Conversely, suppose
\(\cL_s\) is ample.  By \cite[4.7.1]{EGAIII1}, applied to the proper
morphism \(\cX\to\Spec R\) and the closed point \(s\in\Spec R\), there is an
open neighbourhood \(U\ni s\) such that \(\cL\) is relatively ample over
\(U\).  Since \(R\) is local, the only open subset of \(\Spec R\) containing
the closed point is \(\Spec R\) itself; hence \(U=\Spec R\).

\emph{\cref{itm:ch20-sft:cone}.}  By \cref{itm:ch20-sft:nef} and
\cref{itm:ch20-sft:ample} both conditions are read off from the image
\(V\subseteq N^1(\cX_s)\) of \(N^1(\cX/R)\), where they are the nef and
ample conditions of the projective \(k\)-scheme \(\cX_s\).  Nefness is
defined by the closed conditions \((\alpha\cdot C)\geq0\), so the nef cone
is closed and convex.  Because \(\cX\) is projective over \(R\), the
subspace \(V\) contains an ample class; write \(\mathcal A\) for one.  We
use the standard fact that the sum of a nef and an ample class on a
projective scheme over a field is ample.  Consequently, for \(\alpha\in V\)
nef and \(\varepsilon>0\), the class \(\alpha+\varepsilon\mathcal A\) lies in
\(V\) and is ample; letting \(\varepsilon\to0\) exhibits \(\alpha\) as a
limit of ample classes of \(V\), so the nef cone of \(V\) is contained in the
closure of its ample cone.  The reverse inclusion holds because ample
implies nef and the nef cone is closed.  Finally the ample cone of \(V\) is
nonempty, containing \(\mathcal A\).  Only closedness of the nef cone and
the implication ``nef \(+\) ample \(\Rightarrow\) ample'' are used later; we
do not need, and do not assert, openness of the ample cone inside \(V\).
\end{proof}

\begin{remark}[Which cone facts are used]
\label{rem:ch20-cone-facts}
\Cref{lem:ch20-special-fibre-test}\,\cref{itm:ch20-sft:cone} is invoked four
times: in \Cref{prop:ch20-model-max} to make an ample perturbation of a real
class ample; in \Cref{lem:ch20-ample-propagation} and Step~4 of
\Cref{prop:ch20-positive-representative}, both times to add a relatively nef
and a relatively ample class; and in Step~2 of
\Cref{thm:ch20-mixed-hodge} to approximate a relatively nef class by
relatively ample ones.  Every one of these uses only the two facts named at
the end of the proof.  The openness of the ample cone that is used in
Step~1 of \Cref{prop:ch20-positive-representative} is the classical
statement on the smooth projective variety \(X\), not a statement about
\(N^1(\cX/R)\).
\end{remark}

\begin{warning}{Relative nefness is not nefness}
A relatively nef class on \(\cX\) need not be nef on the scheme \(\cX\): the
condition says nothing about curves in the generic fibre, and indeed there
are none that are proper over \(R\).  Symmetrically, a class may restrict to
an ample class on \(X\) and be very far from relatively nef.  Both
directions occur already in \Cref{ex:ch20-p2-blowup-numerics}.  Whenever the
word \emph{nef} appears in \Crefrange{chap:bfj-model-forms}{chap:bfj-continuity}
without qualification it means relatively nef, in the sense of
\Cref{lem:ch20-special-fibre-test}\,\cref{itm:ch20-sft:nef}.
\end{warning}

\begin{lemma}
\label{lem:ch20-pullback-injective}
Let \(\rho:\cY\to\cX\) be a domination of models.  Then
\(\rho^*:N^1(\cX/R)\to N^1(\cY/R)\) is well defined and injective, and it
maps the relatively nef cone into the relatively nef cone.
\end{lemma}

\begin{proof}
Well-definedness and preservation of nefness are
\Cref{prop:nef-pullback}: for an integral vertical curve
\(C\subseteq\cY_s\),

\begin{equation}
\label{eq:ch20-nef-pullback-test}
 (\rho^*\alpha\cdot C)=
 \begin{cases}
  0,&\rho(C)\text{ is a point},\\
  \deg(C/\rho(C))\,(\alpha\cdot\rho(C)),
       &\rho(C)\text{ is a curve},
 \end{cases}
\end{equation}

by the projection formula \Cref{thm:model-projection-formula}.  In
particular \(\rho^*\alpha\equiv_{\cY_s}0\) whenever
\(\alpha\equiv_{\cX_s}0\), so \(\rho^*\) descends to the quotients.

For injectivity, suppose \(\rho^*\alpha\equiv_{\cY_s}0\) and let
\(C\subseteq\cX_s\) be a vertical curve.  Since \(\rho\) is projective,
birational and surjective, \(\rho^{-1}(C)\) contains an irreducible
component \(Z\) dominating \(C\); choose one and let \(d=\dim Z\ge1\), so
that \(\rho|_Z:Z\to C\) is proper and surjective with generic fibre of
dimension \(d-1\).  Fix a relatively very ample line bundle \(\mathcal H\)
on \(\cY\) and consider the one-dimensional cycle class

\[
 \gamma:=c_1(\mathcal H)^{d-1}\cap[Z]\in\operatorname{CH}_1(Z).
\]

Its pushforward is \((\rho|_Z)_*\gamma=m\,[C]\) with

\begin{equation}
\label{eq:ch20-injectivity-cycle}
 m=\deg_{\mathcal H}\bigl(Z_{\eta_C}\bigr)>0,
\end{equation}

the degree of the \((d-1)\)-dimensional generic fibre of \(Z\to C\) with
respect to the very ample \(\mathcal H\); positivity is
positivity of the degree of a nonempty projective scheme under a very ample
embedding.  The projection formula for \(\rho\) now gives

\[
 0=\bigl(\rho^*\alpha\cdot\gamma\bigr)
  =\deg\bigl(c_1(\alpha)\cap(\rho|_Z)_*\gamma\bigr)
  =m\,(\alpha\cdot C),
\]

whence \((\alpha\cdot C)=0\).  As \(C\) was arbitrary,
\(\alpha\equiv_{\cX_s}0\).
\end{proof}

\begin{remark}[Why a cycle class and not a curve]
\label{rem:ch20-cycle-not-curve}
The proof of \Cref{lem:ch20-pullback-injective} works with the cycle class
\(c_1(\mathcal H)^{d-1}\cap[Z]\) rather than with a single curve inside
\(\rho^{-1}(C)\) mapping onto \(C\).  Over an infinite residue field one
could argue with a curve: a general complete intersection of very ample
divisors on an integral projective scheme of dimension \(d\geq2\) is again
integral, by Bertini, and it dominates \(C\).  Over a finite residue field
that form of Bertini is unavailable, and in any case the cycle-class
argument is shorter: only the projection formula and the positivity
\eqref{eq:ch20-injectivity-cycle} of one degree are used, and neither
irreducibility nor genericity is needed.  This is the reason for the
parenthetical remark in the original proof of
\Cref{prop:ch20-semipositive-well-defined}, and it costs nothing.
\end{remark}

\subsection{Two computed numerical spaces}
\label{sec:ch20-n1-examples}

Before continuing it is worth seeing \(N^1(\cX/R)\) in a case in which every
class, curve, and intersection number can be written down.  We will return
to this model repeatedly; it is the smallest surface example in which the
vertical divisors form a space of dimension greater than one.

\begin{example}[The projective plane and one blowup]
\label{ex:ch20-p2-blowup-numerics}
Let \(X=\mathbf P^2_K\), so \(n=2\), and let \(\cX_0=\mathbf P^2_R\) with
special fibre \(F=\mathbf P^2_k\).  Since \(\Pic(\mathbf P^2_R)=\mathbf Z\)
is generated by \(\cO(1)\) and \(\cO(F)=\cO(\divis\pi)\simeq\cO\), the
relative numerical space is

\begin{equation}
\label{eq:ch20-p2-n1}
 N^1(\cX_0/R)=\mathbf R\,h,\qquad h:=c_1(\cO(1)),
 \qquad (h\cdot\ell)=1
\end{equation}

for \(\ell\subset F\) a line.  Now let \(p\in F\) be a closed \(k\)-rational
point and \(\rho:\cX_1=\operatorname{Bl}_p\cX_0\to\cX_0\), with exceptional
divisor \(E\simeq\mathbf P^2_k\) and strict transform \(F_1\) of \(F\).
Because \(F\) is smooth at \(p\) and occurs in \(\cX_{0,s}\) with
multiplicity one,

\begin{equation}
\label{eq:ch20-p2-special-fibre}
 \rho^*F=F_1+E,
 \qquad
 \cX_{1,s}=F_1+E,
 \qquad
 F_1\simeq\operatorname{Bl}_p\mathbf P^2_k,
\end{equation}

and both multiplicities are one.  Write \(e=F_1\cap E\); it is a line in
\(E\simeq\mathbf P^2_k\) and the exceptional curve of \(F_1\), and we write
\(h_1\) for the class in \(F_1\) of the strict transform of a line of
\(\mathbf P^2_k\) missing \(p\).  The three intersection numbers that are
needed everywhere below are

\begin{equation}
\label{eq:ch20-p2-blowup-table}
 (E\cdot\ell_E)=-1,\qquad
 (E\cdot e)=-1,\qquad
 (F_1\cdot e)=+1,
\end{equation}

where \(\ell_E\subset E\) is any line.  Indeed
\(\cO(E)|_E\simeq\cO_{\mathbf P^2}(-1)\), which gives the first two, and
\(\cO(F_1)|_{F_1}=\cO(\rho^*F-E)|_{F_1}=\cO_{F_1}(-e)\), so
\((F_1\cdot e)=-(e^2)_{F_1}=+1\).  As a check,
\((\cX_{1,s}\cdot e)=(F_1\cdot e)+(E\cdot e)=0\), as it must be because
\(\cO(\cX_{1,s})\simeq\cO_{\cX_1}\).

The classes \(\rho^*h\) and \([E]\) are independent in \(N^1(\cX_1/R)\):
testing on \(\ell_E\) gives \((\rho^*h\cdot\ell_E)=0\) and
\(([E]\cdot\ell_E)=-1\), while testing on the strict transform
\(\tilde\ell\subset F_1\) of a line missing \(p\) gives
\((\rho^*h\cdot\tilde\ell)=1\) and \(([E]\cdot\tilde\ell)=(e\cdot h_1)=0\).
Hence

\begin{equation}
\label{eq:ch20-p2-blowup-n1}
 N^1(\cX_1/R)=\mathbf R\,\rho^*h\oplus\mathbf R\,[E],
 \qquad \dim_{\mathbf R}N^1(\cX_1/R)=2,
\end{equation}

and \([F_1]=-[E]\) because \([\cX_{1,s}]=0\).  Both dimensions,
\eqref{eq:ch20-p2-n1} and \eqref{eq:ch20-p2-blowup-n1}, agree with the count
\(\dim N^1(\cX/R)=\dim N^1(X)+\#\{\text{components of }\cX_s\}-1\) that
\Cref{thm:ch20-ddc-exact} will prove in general: here
\(1=1+1-1\) and \(2=1+2-1\).

Finally, \(\mathcal A:=2\rho^*h-[E]\) is relatively ample.  By
\Cref{lem:ch20-special-fibre-test}\,\cref{itm:ch20-sft:ample} it suffices to
check ampleness on \(\cX_{1,s}=F_1\cup E\), hence on each component:
\(\mathcal A|_E=\cO_{\mathbf P^2}(1)\), and
\(\mathcal A|_{F_1}=2h_1-e\), which is ample on
\(\operatorname{Bl}_p\mathbf P^2_k\) by the Nakai--Moishezon criterion, since
\((2h_1-e)^2=3>0\), \((2h_1-e)\cdot h_1=2>0\), \((2h_1-e)\cdot e=1>0\), and
every irreducible curve on \(\operatorname{Bl}_p\mathbf P^2_k\) has class
\(ah_1-be\) with \(a>0\), \(a\ge b\), or is \(e\) itself.
\end{example}

\section{Closed model forms}
\label{sec:ch20-model-forms}

\subsection{The direct system}
\label{sec:ch20-direct-system}

Let \(\mathfrak M(X)\) denote the collection of models of \(X\) in the sense
of \eqref{eq:ch20-standing}, preordered by domination: \(\cY\geq\cX\) if
there is a (necessarily unique) domination \(\rho_{\cY\cX}:\cY\to\cX\).
Uniqueness holds because two dominations agree on the dense open subscheme
\(X\subset\cY\) and \(\cX\) is separated; consequently
\(\rho_{\mathcal Z\cX}=\rho_{\cY\cX}\circ\rho_{\mathcal Z\cY}\) whenever
\(\mathcal Z\geq\cY\geq\cX\), and \(\mathfrak M(X)\) is a genuine directed
preordered index category rather than merely a set with a relation.
Directedness is \Cref{prop:common-domination}: two models are dominated by
the closure of the graph of the birational map between them, taken inside
their fibre product and normalized.

By \Cref{lem:ch20-pullback-injective} the assignment
\(\cX\mapsto N^1(\cX/R)\), \(\rho\mapsto\rho^*\), is a functor from
\(\mathfrak M(X)\) to finite-dimensional real vector spaces with injective
transition maps.  Its direct limit is therefore the union of an increasing
family of finite-dimensional spaces, and every element of the limit is
represented on some model.

\begin{definition}[Closed model forms]
\index{closed model form}\index{direct limit!over models}
\label{def:ch20-closed-model-form}
The space of \emph{closed model \((1,1)\)-forms} is

\begin{equation}
\label{eq:ch20-model-forms}
 \mathcal Z^{1,1}(X)
   :=\varinjlim_{\mathcal X}N^1(\mathcal X/R).
\end{equation}

A form \(\theta\) is \emph{determined} on \(\mathcal X\) if it is represented
by a class \(\theta_{\mathcal X}\in N^1(\mathcal X/R)\).  Its generic
numerical class is denoted

\begin{equation}
\label{eq:ch20-generic-class}
                 \{\theta\}\in N^1(X).
\end{equation}
\end{definition}

Concretely, an element of \(\mathcal Z^{1,1}(X)\) is an equivalence class of pairs
\((\cX,\theta_{\cX})\), where \((\cX,\theta_{\cX})\sim(\cY,\theta_{\cY})\)
if there is a model \(\mathcal Z\) dominating both with
\(\rho_{\mathcal Z\cX}^*\theta_{\cX}=\rho_{\mathcal Z\cY}^*\theta_{\cY}\).  The map
\(\{\cdot\}\) of \eqref{eq:ch20-generic-class} is induced by restriction to
the generic fibre, \(N^1(\cX/R)\to N^1(X)\), which is compatible with
pullback because \(\rho\) restricts to the identity on \(X\).

Thus equality of forms never means equality on every model.  It means that
the two representatives become equal after pullback to one common model.
The braces in \eqref{eq:ch20-generic-class} are called the de Rham class in
the BFJ literature, although only numerical equivalence is used in this
chapter.

The next proposition is the statement that makes the limit usable: the
existential quantifier in the definition of \(\sim\) may be replaced by a
universal one, so that a single computation on a single model decides
equality.

\begin{proposition}
\label{prop:ch20-two-model-comparison}
Let \(\theta_{\cX}\in N^1(\cX/R)\) and \(\theta_{\cY}\in N^1(\cY/R)\), and
let \(\mathcal Z\) be any model dominating both.  The following are equivalent.
\begin{enumerate}
\item\label{itm:ch20-tmc:some} \((\cX,\theta_{\cX})\) and
 \((\cY,\theta_{\cY})\) define the same element of \(\mathcal Z^{1,1}(X)\);
\item\label{itm:ch20-tmc:this}
 \(\rho_{\mathcal Z\cX}^*\theta_{\cX}=\rho_{\mathcal Z\cY}^*\theta_{\cY}\) in
 \(N^1(\mathcal Z/R)\);
\item\label{itm:ch20-tmc:all}
 \(\rho_{\mathcal W\cX}^*\theta_{\cX}=\rho_{\mathcal W\cY}^*\theta_{\cY}\) for
 \emph{every} model \(\mathcal W\) dominating \(\cX\) and \(\cY\);
\item\label{itm:ch20-tmc:curves}
 \((\rho_{\mathcal Z\cX}^*\theta_{\cX}\cdot C)
   =(\rho_{\mathcal Z\cY}^*\theta_{\cY}\cdot C)\) for every vertical curve
 \(C\subseteq\mathcal Z_s\).
\end{enumerate}
\end{proposition}

\begin{proof}
\cref{itm:ch20-tmc:all}\(\Rightarrow\)\cref{itm:ch20-tmc:this}\(\Rightarrow\)\cref{itm:ch20-tmc:some}
is trivial, and \cref{itm:ch20-tmc:this}\(\Leftrightarrow\)\cref{itm:ch20-tmc:curves}
is \Cref{def:ch20-relative-numerical}.  It remains to prove
\cref{itm:ch20-tmc:some}\(\Rightarrow\)\cref{itm:ch20-tmc:all}.  Suppose the
two pairs agree on some model \(\mathcal Z'\) dominating both, and let \(\mathcal W\) be
an arbitrary model dominating \(\cX\) and \(\cY\).  By
\Cref{prop:common-domination} choose \(\mathcal V\) dominating \(\mathcal W\) and
\(\mathcal Z'\).  All the dominations in sight are compatible by the uniqueness
noted above, so

\[
 \rho_{\mathcal V\mathcal W}^*\rho_{\mathcal W\cX}^*\theta_{\cX}
 =\rho_{\mathcal V\cX}^*\theta_{\cX}
 =\rho_{\mathcal V\mathcal Z'}^*\rho_{\mathcal Z'\cX}^*\theta_{\cX}
 =\rho_{\mathcal V\mathcal Z'}^*\rho_{\mathcal Z'\cY}^*\theta_{\cY}
 =\rho_{\mathcal V\mathcal W}^*\rho_{\mathcal W\cY}^*\theta_{\cY}.
\]

Since \(\rho_{\mathcal V\mathcal W}^*\) is injective
(\Cref{lem:ch20-pullback-injective}), the two classes on \(\mathcal W\) coincide.
\end{proof}

\begin{corollary}
\label{cor:ch20-limit-structure}
\(\mathcal Z^{1,1}(X)\) is a real vector space; every element is determined on some
model; any finite family of elements is determined on a common model; and
for each \(\cX\) the natural map \(N^1(\cX/R)\to\mathcal Z^{1,1}(X)\) is injective.
Moreover \(\{\cdot\}:\mathcal Z^{1,1}(X)\to N^1(X)\) is linear and surjective.
\end{corollary}

\begin{proof}
The first four assertions are the standard properties of a direct limit over
a directed index category with injective transition maps, using
\Cref{prop:ch20-two-model-comparison} for injectivity of
\(N^1(\cX/R)\to\mathcal Z^{1,1}(X)\).  Linearity of \(\{\cdot\}\) is clear.  For
surjectivity, let \(\alpha\in N^1(X)\) and write it as a finite real
combination of classes of line bundles \(L_j\) on \(X\).  By
\Cref{prop:line-bundle-extension} each \(L_j\) has a multiple that extends
to a line bundle on some model; pass to a common model by
\Cref{prop:common-domination} and take the corresponding real combination.
\end{proof}

\begin{warning}{A closed model form is not a form}
The notation \(\mathcal Z^{1,1}\) and the word \emph{form} are borrowed from
complex geometry and are convenient, but nothing in
\Cref{def:ch20-closed-model-form} is a differential form, and no element of
\(\mathcal Z^{1,1}(X)\) is a function on \(X^{\An}\).  A closed model form is a
compatible system of numerical classes.  The only sense in which it can be
evaluated is through intersection numbers
(\Cref{sec:ch20-intersection}) and through the measures of
\Cref{chap:bfj-ma}.  The bridge to genuine differential objects is
\Cref{thm:ch19-cld-bfj-comparison}, and it exists only on the model and
continuous loci described there.
\end{warning}

\subsection{Semipositive and positive forms}
\label{sec:ch20-semipositive-def}

\begin{definition}[Semipositivity of a closed model form]
\index{semipositivity!of a closed model form}
\label{def:ch20-semipositive}
A closed model form \(\theta\) is \emph{semipositive} if it has a
determination \(\theta_{\mathcal X}\) that is relatively nef.  It is
\emph{\(\mathcal X\)-positive} if it is determined on \(\mathcal X\) by a
relatively ample class.
\end{definition}

The distinction matters.  Nefness is preserved and detected by a vertical
modification, whereas the pullback of an ample class is generally only nef.
Accordingly \(\mathcal X\)-positivity is a statement about a specific model
and is \emph{not} a property of \(\theta\) alone, whereas semipositivity is
a property of \(\theta\); \Cref{prop:ch20-semipositive-well-defined} is the
proof of the latter, and \Cref{ex:ch20-nef-not-ample} shows that the former
really does depend on the model.

\begin{proposition}\label{prop:ch20-semipositive-well-defined}
Semipositivity is independent of the determination.  The semipositive forms
form a convex cone, and pullback by a morphism of models preserves
semipositivity.
\end{proposition}

\begin{proof}
Let \(\theta_{\cX}\) be a determination and \(\rho:\cY\to\cX\) a
domination.  By \eqref{eq:ch20-nef-pullback-test}, relative nefness of
\(\theta_{\cX}\) implies relative nefness of \(\rho^*\theta_{\cX}\); this is
the easy half and is \Cref{prop:nef-pullback}.

Conversely, suppose \(\rho^*\theta_{\cX}\) is relatively nef and let
\(C\subseteq\cX_s\) be a vertical curve.  Exactly the argument of
\Cref{lem:ch20-pullback-injective} applies: choose an irreducible component
\(Z\) of \(\rho^{-1}(C)\) dominating \(C\), let
\(\gamma=c_1(\mathcal H)^{\dim Z-1}\cap[Z]\) for a relatively very ample
\(\mathcal H\), and use the projection formula together with
\eqref{eq:ch20-injectivity-cycle}:

\[
 0\leq(\rho^*\theta_{\cX}\cdot\gamma)=m\,(\theta_{\cX}\cdot C),
 \qquad m>0 .
\]

The first inequality holds because \(\gamma\) is a nonnegative combination of
classes of curves on \(\cY_s\)---indeed \(\mathcal H\) is very ample, so
\(c_1(\mathcal H)^{\dim Z-1}\cap[Z]\) is represented by an effective
one-cycle---and \(\rho^*\theta_{\cX}\) is relatively nef.  Hence
\((\theta_{\cX}\cdot C)\geq0\), and \(\theta_{\cX}\) is relatively nef.
Working with a cycle class rather than a single curve is what makes the
argument valid over a finite residue field
(\Cref{rem:ch20-cycle-not-curve}).

Now let \(\theta_{\cX}\) and \(\theta_{\cY}\) be two determinations of the
same \(\theta\) and let \(\mathcal Z\) dominate both.  By
\Cref{prop:ch20-two-model-comparison} their pullbacks to \(\mathcal Z\) agree, and
by the two paragraphs above each of them is relatively nef if and only if
that common pullback is.  Hence relative nefness of one determination is
equivalent to relative nefness of the other, which is the asserted
independence.

Convexity: if \(\theta,\theta'\) are semipositive and \(s,t\geq0\), choose a
common model on which both are determined by relatively nef classes
(\Cref{cor:ch20-limit-structure}); then
\(s\theta_{\cX}+t\theta'_{\cX}\) has nonnegative degree on every vertical
curve, so \(s\theta+t\theta'\) is semipositive.

Stability under pullback: let \(\rho:\cY\to\cX\) be an arbitrary morphism of
models over \(R\), not necessarily birational and not necessarily with the
same generic fibre, and let \(\theta_{\cX}\) be relatively nef.  A vertical
curve \(C\subseteq\cY_s\) has \(\rho(C)\subseteq\cX_s\) either a point or a
vertical curve, since \(\rho\) is proper and preserves special fibres; in
the first case \((\rho^*\theta_{\cX}\cdot C)=0\) and in the second it equals
\(\deg(C/\rho(C))\,(\theta_{\cX}\cdot\rho(C))\geq0\).  This is
\eqref{eq:ch20-nef-pullback-test} again, and it proves that
\(\rho^*\theta_{\cX}\) is relatively nef.
\end{proof}

\begin{corollary}
\label{cor:ch20-pullback-varieties}
Let \(f:X'\to X\) be a morphism of smooth projective \(K\)-varieties and let
\(\theta\in\mathcal Z^{1,1}(X)\) be semipositive.  Then there is a
well-defined semipositive \(f^*\theta\in\mathcal Z^{1,1}(X')\), and
\(f^*(\theta+\ddc\varphi)=f^*\theta+\ddc(f^*\varphi)\) for every model
function \(\varphi\) on \(X^{\An}\).
\end{corollary}

\begin{proof}
Let \(\theta\) be determined on \(\cX\).  Choose any model \(\cX'_0\) of
\(X'\) and let \(\cX'\) be the flat closure of the graph of
\((\mathrm{id},f):X'\to\cX'_0\times_R\cX\), a model of \(X'\) equipped with a
morphism \(\rho:\cX'\to\cX\) extending \(f\) (this is the graph-closure
construction of \Cref{prop:common-domination}, applied to a morphism rather
than to a birational map).  Set \(f^*\theta:=[\rho^*\theta_{\cX}]\), which is
semipositive by the previous proof.  Independence of the choices follows by
comparing two such constructions on a common domination, using
\Cref{prop:ch20-two-model-comparison}.  The compatibility with \(\ddc\) is
the identity \(\rho^*[D]=[\rho^*D]\) together with
\(\varphi_{\rho^*D}=\varphi_D\circ f^{\An}\).
\end{proof}

\begin{corollary}
\label{cor:ch20-semipositive-generic-nef}
If \(\theta\in\mathcal Z^{1,1}(X)\) is semipositive, then \(\{\theta\}\) is nef on
\(X\).
\end{corollary}

\begin{proof}
Let \(\theta_{\cX}\) be a relatively nef determination and let \(C\subset X\)
be an integral curve.  Let \(\overline C\subseteq\cX\) be its flat closure,
an integral two-dimensional scheme, flat and proper over \(R\) with generic
fibre \(C\) (\Cref{lem:flat-closure-generic-fibre}), and let
\(\nu:\mathcal C\to\overline C\) be its normalization, so that
\(\mathcal C\) is a normal proper flat model of the normalization
\(\widetilde C\) of \(C\).  Write \(\mu:\widetilde C\to C\) for the induced
finite birational morphism on generic fibres, of degree one.

By the projection formula on \(C\),
\((\{\theta\}\cdot C)_X=\deg_{\widetilde C}\mu^*(\{\theta\}|_C)\).  By
conservation of degree in the flat proper family
\(\mathcal C\to\Spec R\), that is by
\Cref{cor:special-fibre-relation} applied to \(\mathcal C\) with the single
line bundle \(\nu^*(\theta_{\cX}|_{\overline C})\),

\[
 (\{\theta\}\cdot C)_X
 =\sum_{C'}b_{C'}\,\bigl(\nu^*\theta_{\cX}\cdot C'\bigr),
\]

the sum running over the irreducible components \(C'\) of the special fibre
of \(\mathcal C\), with their multiplicities \(b_{C'}>0\).  Each \(C'\) maps
finitely onto a vertical curve of \(\cX\) or to a point; in the first case
\((\nu^*\theta_{\cX}\cdot C')=\deg(C'/\nu(C'))(\theta_{\cX}\cdot\nu(C'))\geq0\)
by the projection formula \eqref{eq:ch20-nef-pullback-test}, and in the
second case it vanishes.  Every term is therefore nonnegative.
\end{proof}

The converse of \Cref{cor:ch20-semipositive-generic-nef} is false, and
badly so: by \Cref{thm:ch20-vertical-nef-rigidity} below, the forms
\(\ddc\varphi\) with \(\varphi\) a nonconstant model function all have
generic class \(0\), which is nef, and none of them is semipositive.  What
\emph{is} always available for an ample generic class is a positive
representative after adding a model potential, and this is the statement
that the rest of the theory uses.

That statement is \Cref{prop:ch20-positive-representative}.  Its proof needs
the exact sequence of \Cref{sec:ch20-exact}---specifically, the fact that a
class with numerically trivial generic restriction is the class of a
vertical divisor---so it is deferred to
\Cref{sec:ch20-positive-representative}, immediately after that sequence has
been established.  Nothing between here and there uses it.

\begin{example}[Nef but not \(\mathcal X\)-positive]
\label{ex:ch20-nef-not-ample}
In \Cref{ex:ch20-p2-blowup-numerics} take \(\theta=\rho^*h\) on
\(\mathcal X_1=\operatorname{Bl}_p\mathbf P^2_R\).  It is semipositive: its
determination on \(\mathcal X_0\) is \(h\), which is relatively ample there,
so \(\rho^*h\) is relatively nef.  It is not \(\mathcal X_1\)-positive,
because \((\rho^*h\cdot\ell_E)=0\) for a line \(\ell_E\subseteq E\).  Thus
\(\mathcal X\)-positivity genuinely depends on \(\mathcal X\), while
semipositivity does not.  The repair supplied by
\Cref{prop:ch20-positive-representative} is visible here: the classes on
\(\mathcal X_1\) with generic class \(h\) are \(\rho^*h+a[E]\), with degrees
\(-a\) on a line of \(E\) and \(1+a\) on the strict transform of a line
through \(p\), so such a class is relatively nef exactly for
\(-1\leq a\leq0\) and relatively ample exactly for \(-1<a<0\).  Taking
\(a=-\tfrac12\), that is \(\psi=-\tfrac12\varphi_E\), makes
\(\theta+\ddc\psi\) relatively ample; taking \(a=-1\) overshoots, since
\((\rho^*h-[E])|_{F_1}=h_1-e\) is nef but not ample.  The example is a
reminder that the cofinal conclusion of
\Cref{prop:ch20-positive-representative} cannot be strengthened to ``every
model'', and that the correcting potential is not unique.
\end{example}

\section{Model functions and the model differential}
\label{sec:ch20-ddc}

\subsection{The space of model functions}
\label{sec:ch20-model-function-space}

Let \(\Div_0(\mathcal X)_{\mathbf R}\) be the real vector space of vertical
Cartier divisors on a model \(\mathcal X\), that is, of \(\mathbf R\)-Cartier
divisors supported on \(\mathcal X_s\).  If \(\mathcal X\) is normal and
\(\mathcal X_s=\sum_{i\in I}b_iE_i\) is its special fibre with its
multiplicities, then \(\Div_0(\mathcal X)_{\mathbf R}\) is contained in
\(\bigoplus_{i\in I}\mathbf R\,E_i\), with equality when \(\mathcal X\) is
regular, since then each \(E_i\) is Cartier.  Pullback along a domination
sends \(\Div_0(\mathcal X)_{\mathbf R}\) into
\(\Div_0(\mathcal Y)_{\mathbf R}\) and is injective, because it is injective
already on the associated model functions
(\Cref{prop:metrics-divisor-model-function}\,(2)).  Pullback over all models
therefore gives the space of real model functions

\begin{equation}
\label{eq:ch20-model-functions}
 \mathcal D(X)_{\mathbf R}
   =\varinjlim_{\mathcal X}\Div_0(\mathcal X)_{\mathbf R},
\end{equation}

embedded in \(C^0(X^{\An})\) by \(D\mapsto\varphi_D\).  Rational model
functions form \(\mathcal D(X)\).  We use real model functions when taking
convex combinations and rational ones when invoking line bundles or base
ideals.

Three facts about this space, all proved in \Cref{part:foundations},
will be used without further comment.  First, \(\mathcal D(X)_{\mathbf Q}\)
is uniformly dense in \(C^0(X^{\An})\)
(\Cref{prop:metrics-model-function-density}); this is what makes model
functions an adequate test class for Radon measures in
\Cref{chap:bfj-ma}.  Second, a model function determined on an SNC model
\(\mathcal X\) is affine on every face of the dual complex
\(\Delta_{\mathcal X}\) and satisfies
\(\varphi_D\circ r_{\mathcal X}=\varphi_D\)
(\Cref{lem:model-function-affine}); in particular it is determined by its
values \(\varphi_D(x_{E_i})=a_i/b_i\) at the divisorial points of
\(\mathcal X\).  Third, the correspondence
\(D\mapsto\varphi_D|_{\Delta_{\mathcal X}}\) is an isomorphism onto the
rational affine functions on the faces of \(\Delta_{\mathcal X}\)
(\Cref{prop:vertical-divisor-affine-iso}).  Together these say that
\(\mathcal D(X)_{\mathbf R}\) is the space of \emph{piecewise affine
functions on the inverse system of skeleta}, and that the model-theoretic
and skeletal descriptions of a potential are interchangeable.

\subsection{The model differential}
\label{sec:ch20-ddc-definition}

\begin{definition}[The model differential]
\index{ddc operator@\(\ddc\) operator!on model functions}\index{model function}
\label{def:ch20-ddc}
For \(\varphi_D\in\mathcal D(X)_{\mathbf R}\), define

\begin{equation}
\label{eq:ch20-ddc}
                         \ddc\varphi_D=[D]
                         \in\mathcal Z^{1,1}(X).
\end{equation}

If \(\theta\in\mathcal Z^{1,1}(X)\), the model function \(\varphi_D\) is
\emph{\(\theta\)-psh at model level} if
\(\theta+\ddc\varphi_D\) is semipositive.
\end{definition}

\begin{lemma}
\label{lem:ch20-ddc-well-defined}
The prescription \eqref{eq:ch20-ddc} depends only on the function
\(\varphi_D\) and not on the pair \((\mathcal X,D)\) representing it, and it
is \(\mathbf R\)-linear.  Moreover \(\ddc\) commutes with pullback in the
sense that, for a domination \(\rho:\mathcal Y\to\mathcal X\) and
\(D\in\Div_0(\mathcal X)_{\mathbf R}\), the class of \(\rho^*D\) on
\(\mathcal Y\) represents the same closed model form as the class of \(D\)
on \(\mathcal X\).
\end{lemma}

\begin{proof}
Let \((\mathcal X,D)\) and \((\mathcal X',D')\) satisfy
\(\varphi_D=\varphi_{D'}\) as functions on \(X^{\An}\), and let
\(\rho:\mathcal Y\to\mathcal X\), \(\rho':\mathcal Y\to\mathcal X'\) be
dominations from a common model (\Cref{prop:common-domination}).  Pullback
of a Cartier divisor pulls back its local equations, hence its model
function: \(\varphi_{\rho^*D}=\varphi_D\circ\iota=\varphi_D\), where we use
that \(\rho\) induces the identity on \(X^{\An}\).  Therefore
\(\varphi_{\rho^*D}=\varphi_{\rho'^*D'}\) on \(X^{\An}\), and by injectivity
of \(D\mapsto\varphi_D\) on a fixed normal model
(\Cref{prop:divisor-model-function-injective-mod-fibre}) we get
\(\rho^*D=\rho'^*D'\) as divisors on \(\mathcal Y\).  A fortiori their
numerical classes agree, so by \Cref{prop:ch20-two-model-comparison} the two
pairs define the same element of \(\mathcal Z^{1,1}(X)\).  Linearity is
inherited from linearity of \(D\mapsto\varphi_D\) and of
\(D\mapsto[\cO(D)]\).  The last assertion was proved along the way.
\end{proof}

Since \(c\mathcal X_s=\operatorname{div}(\pi^c)\) as
a real Cartier divisor and \(\varphi_{c\mathcal X_s}=c\), one has

\begin{equation}
\label{eq:ch20-ddc-constant}
                            \ddc c=0.
\end{equation}

\Cref{thm:ch20-ddc-exact} will show that constants are the \emph{only}
model functions killed by \(\ddc\).

Equation \eqref{eq:ch20-ddc-constant} does not say that the special-fibre
cycle has zero relative degree.  If
\(\mathcal L_1,\ldots,\mathcal L_n\) extend \(L_1,\ldots,L_n\), then

\begin{equation}
\label{eq:ch20-special-fibre-degree}
 (\mathcal X_s\cdot\mathcal L_1\cdots\mathcal L_n)
                  =(L_1\cdots L_n)_X.
\end{equation}

The left side treats \(\mathcal X_s\) as a vertical cycle; the zero in
\eqref{eq:ch20-ddc-constant} treats its associated line bundle as a
numerical class.  Keeping these two roles separate prevents a later
mass-zero sign error.

\begin{warning}{Two roles of the special fibre, and one recurring error}
The identity \eqref{eq:ch20-special-fibre-degree} is
\Cref{cor:special-fibre-relation}: it is conservation of degree in the flat
family \(\mathcal X\to\Spec R\), and its right-hand side is positive as soon
as the \(L_i\) are ample.  The identity \eqref{eq:ch20-ddc-constant} says
that \(\cO_{\mathcal X}(\mathcal X_s)\simeq\cO_{\mathcal X}\), so that the
class of \(\mathcal X_s\) in \(N^1(\mathcal X/R)\) vanishes.  Both are true,
and confusing them produces a total mass of zero in
\Cref{chap:bfj-ma}.  The distinction is that
\eqref{eq:ch20-special-fibre-degree} intersects the \emph{cycle}
\(\sum_ib_i[E_i]\) with \(n\) line bundles, obtaining a number, while
\eqref{eq:ch20-ddc-constant} intersects the \emph{line bundle}
\(\cO(\mathcal X_s)\) with \(n-1\) line bundles and a curve, obtaining
zero.  The safe rule is \eqref{eq:ch20-sign}: \(\varphi_{\mathcal X_s}=1\),
and integrating the constant \(1\) against a Chambert--Loir measure returns
its total mass \((L_1\cdots L_n)\) by
\Cref{prop:metrics-cl-model-function-integral}.
\end{warning}

\subsection{Normalization audit}
\label{sec:ch20-normalization-audit}

Before any computation we check \eqref{eq:ch20-ddc} against the unit tests
of \Cref{sec:appJ-unit-tests}.  The point of the exercise is that
\(\ddc\) has now been defined twice in this book---once here, on model
functions, with values in numerical classes, and once in
\Cref{app:normalizations} as the complex operator
\eqref{eq:appJ-ddc}---and a third operator \(\dprimeddprime\) exists in
\Cref{part:cld}.  They are related but not equal, and the relation is a
theorem, not a convention.

\begin{enumerate}
\item \emph{Constant test.}  \(\ddc(\varphi+c)=\ddc\varphi\) by
 \eqref{eq:ch20-ddc-constant}.  Consequently
 \(\theta+\ddc\varphi\) is unchanged, and so are all the Monge--Amp\`ere
 measures built from it in \Cref{chap:bfj-ma}.  Energies change by
 \(c\,V\) with \(V=(\{\theta\}^n)\), exactly as
 \eqref{eq:appJ-energy-translation} prescribes.
\item \emph{Special-fibre test.}  \(D=\mathcal X_s\) gives
 \(\varphi_D=+1\), by \eqref{eq:ch20-sign} and
 \eqref{eq:appJ-special-fibre-test}.  A source assigning \(-1\) to the
 special fibre is using the opposite metric convention and must be
 translated before its formulae are imported.
\item \emph{Effectivity test.}  \(D\geq0\) if and only if
 \(\varphi_D\geq0\), by
 \Cref{prop:metrics-divisor-model-function}\,(4).  Thus an ideal
 \(\mathfrak a\subseteq\cO_{\mathcal X}\) has
 \(\log|\mathfrak a|\leq0\); see \eqref{eq:ch20-ideal-divisor}.
\item \emph{Principal-divisor test.}  If \(D=\divis(g)\) is principal on
 \(\mathcal X\) with \(g\in K(X)^\times\), then \([D]=0\) in
 \(N^1(\mathcal X/R)\), so \(\ddc\varphi_D=0\); by
 \Cref{thm:ch20-ddc-exact} this forces \(\varphi_D\) to be constant.  This
 is consistent with \Cref{prop:metrics-divisor-model-function}\,(5), which
 warns that \(\varphi_{\divis(g)}=-\log|g|\) need not vanish
 identically---but it must be constant when \(\divis(g)\) is vertical,
 because then \(g\) is a unit on \(X\) and \(X\) is proper and connected, so
 \(g\in K^\times\) and \(-\log|g|=v_K(g)\).
\item \emph{Total-mass test.}  Integrating \(\varphi_D\) against a model
 Chambert--Loir measure returns the vertical intersection number
 \eqref{eq:metrics-cl-pairing}, and taking \(D=\mathcal X_s\) returns
 \((L_1\cdots L_n)\).
\end{enumerate}

\begin{warning}{\(\ddc\) here is not \(\dprimeddprime\)}
In \Cref{app:normalizations}, \eqref{eq:appJ-valuation-pl} records
\(\dprimeddprime[v(f)]=-\delta_{\divis(f)}\) for the
Chambert--Loir--Ducros operator, whereas
\Cref{def:ch20-ddc} gives \(\ddc\varphi_D=+[D]\) with
\(\varphi_D=v(f_D)\).  There is no contradiction and no sign error: the two
operators act on different objects and land in different targets.  The CLD
operator is a genuine local differential operator applied to a function on
\(X^{\An}\), and it produces a current on \(X^{\An}\); the operator of
\Cref{def:ch20-ddc} is the passage from a vertical divisor to its
\emph{relative numerical class on a model}, and it produces an element of
\(\mathcal Z^{1,1}(X)\), which is not a current.  The divisor
\(\divis(f)\) appearing in \eqref{eq:appJ-valuation-pl} is horizontal, being
the divisor of a function on the analytic space; the divisor \(D\) appearing
in \eqref{eq:ch20-ddc} is vertical.  The precise statement relating the two
is \Cref{thm:ch19-cld-bfj-comparison}: for a model function \(\varphi\) and
a model metric \(\|\cdot\|_0\) with closed model form \(\theta\), the CLD
Chern current of \(\|\cdot\|_\varphi=\rme^{-\varphi}\|\cdot\|_0\) is
positive exactly when \(\theta+\ddc\varphi\) is semipositive, and the two
top products agree as measures.  Outside the model and continuous loci that
theorem asserts nothing, and neither do we.
\end{warning}

\subsection{The model differential on a curve}
\label{sec:ch20-ddc-curve}

The case \(n=1\) is the one in which \(\ddc\) can be written as an explicit
finite matrix, and it is worth recording once and for all, because the
higher-dimensional formulae in \Cref{chap:bfj-compactness} are modelled on
it.

\begin{proposition}
\label{prop:ch20-curve-laplacian}
Let \(n=1\), let \(\mathcal X\) be a regular model of the smooth projective
curve \(X\), and write \(\mathcal X_s=\sum_{i\in I}b_iE_i\).  Let
\(D=\sum_ia_iE_i\) be a vertical \(\mathbf R\)-divisor, with associated
vertex values \(u_i=\varphi_D(x_{E_i})=a_i/b_i\).  Then, for every
\(i\in I\),

\begin{equation}
\label{eq:ch20-curve-laplacian}
 (D\cdot E_i)=\sum_{j\neq i}b_j(E_j\cdot E_i)\,(u_j-u_i),
 \qquad
 \sum_{i}b_i\,(D\cdot E_i)=0 ,
\end{equation}

and the weights \(b_j(E_j\cdot E_i)\), \(j\ne i\), are nonnegative.
Consequently the map \(u\mapsto\bigl((D\cdot E_i)\bigr)_i\) is the graph
Laplacian of the weighted incidence graph of \(\mathcal X_s\).  If moreover
\(\mathcal X\) is semistable, so that all \(b_i=1\), all singular points of
\(\mathcal X_s\) are nodes with local equation \(xy=\pi\), and all edges of
\(\Delta_{\mathcal X}\) have length one, then

\begin{equation}
\label{eq:ch20-curve-laplacian-graph}
 \sum_{i}(D\cdot E_i)\,\delta_{x_{E_i}}
 =\Delta_{\Delta_{\mathcal X}}\varphi_D
\end{equation}

with the sign convention of \eqref{eq:graph-laplacian}: a model function with
a convex corner at a vertex has positive Laplacian there.
\end{proposition}

\begin{proof}
Write \(D=\sum_jb_ju_jE_j\).  Because
\(\cO_{\mathcal X}(\mathcal X_s)\simeq\cO_{\mathcal X}\) we have
\(\sum_jb_j(E_j\cdot E_i)=(\mathcal X_s\cdot E_i)=0\) for each \(i\), so
subtracting \(u_i\) times this identity from
\((D\cdot E_i)=\sum_jb_ju_j(E_j\cdot E_i)\) leaves
\(\sum_{j}b_j(E_j\cdot E_i)(u_j-u_i)\), and the term \(j=i\) vanishes.  This
is the first identity in \eqref{eq:ch20-curve-laplacian}.  The second is
\((D\cdot\mathcal X_s)=0\) again, now with the roles of \(D\) and
\(\mathcal X_s\) exchanged, using symmetry of the intersection pairing on
the regular surface \(\mathcal X\).  For \(j\neq i\), \(E_i\) and \(E_j\)
are distinct prime divisors on the regular two-dimensional scheme
\(\mathcal X\), so their intersection is a finite scheme and
\((E_j\cdot E_i)\geq0\); \(b_j>0\) by definition.

For \eqref{eq:ch20-curve-laplacian-graph}, assume \(\mathcal X\) semistable.
The edges of \(\Delta_{\mathcal X}\) are in bijection with the nodes of
\(\mathcal X_s\), an edge joining \(x_{E_i}\) to \(x_{E_j}\) for each point
of \(E_i\cap E_j\); by \Cref{prop:ch08-annulus-skeleton}\,(3) each has length
\(1\), since the local equation is \(xy=\pi\) and \(v_K(\pi)=1\).  The number
of edges joining \(x_{E_i}\) to \(x_{E_j}\) is therefore \((E_i\cdot E_j)\).
By \Cref{lem:model-function-affine} the function \(\varphi_D\) is affine on
each edge with endpoint values \(u_i,u_j\), so its outgoing slope at
\(x_{E_i}\) along that edge is \(u_j-u_i\).  Summing the outgoing slopes at
\(x_{E_i}\) over all edges emanating from it gives
\(\sum_{j\neq i}(E_i\cdot E_j)(u_j-u_i)\), which by the first identity in
\eqref{eq:ch20-curve-laplacian} is \((D\cdot E_i)\).  Comparing with
\eqref{eq:graph-laplacian}, whose atom at \(x_{E_i}\) is exactly the sum of
outgoing slopes, gives \eqref{eq:ch20-curve-laplacian-graph}.  A model
function determined on \(\mathcal X\) has no slope breaks in the interior of
an edge, so no other atoms occur.
\end{proof}

\begin{remark}[Multiplicities and edge lengths]
\label{rem:ch20-multiplicities-lengths}
Identity \eqref{eq:ch20-curve-laplacian-graph} was stated only in the
semistable case, and the restriction is not cosmetic.  For a general regular
model the edge of \(\Delta_{\mathcal X}\) attached to a point of
\(E_i\cap E_j\) has length \(1/(b_ib_j)\) in the integral affine structure
of \Cref{def:integral-affine}, and the divisorial point is
\(x_{E_i}=b_i^{-1}\ord_{E_i}\).  Both conventions must be applied together,
as \Cref{rem:divisorial-normalization} and the second remark of
\Cref{sec:appJ-valuations} insist.  The first identity in
\eqref{eq:ch20-curve-laplacian} is valid regardless; it is only the
identification of its right-hand side with the metric-graph Laplacian
\eqref{eq:graph-laplacian} that carries the normalization.
\end{remark}

\begin{example}[A model form on the running conic, carried to numbers]
\label{ex:ch20-conic-ddc}
Take the running family \eqref{eq:running-conic},
\(\mathcal C=\{XY=\pi Z^2\}\subset\mathbf P^2_R\), whose generic fibre is a
smooth conic \(C\) of genus zero and whose special fibre is
\(C_0\cup C_1\), two lines meeting transversally at one point, both with
multiplicity one.  Let
\(\theta=c_1(\cL)\) with \(\cL=\cO_{\mathbf P^2_R}(1)|_{\mathcal C}\), so
that \(d_0=(\theta\cdot C_0)=1\) and \(d_1=(\theta\cdot C_1)=1\), and
\(\{\theta\}=c_1(L)\) has degree \(2\) on \(C\).  The dual graph is a single
edge of length one with vertices \(v_0,v_1\), and the intersection matrix is

\begin{equation}
\label{eq:ch20-conic-matrix}
 \bigl((C_i\cdot C_j)\bigr)_{i,j=0,1}
 =\begin{pmatrix}-1&1\\1&-1\end{pmatrix},
\end{equation}

since \(((C_0+C_1)\cdot C_i)=0\) and \((C_0\cdot C_1)=1\).

Let \(\varphi\) be the model function with \(\varphi(v_0)=A\),
\(\varphi(v_1)=B\); thus \(D=AC_0+BC_1\) and \(\varphi=\varphi_D\).  By
\eqref{eq:ch20-curve-laplacian},

\[
 (D\cdot C_0)=B-A,
 \qquad
 (D\cdot C_1)=A-B,
\]

so \(\ddc\varphi\) is the closed model form whose degrees on the two
components are \((B-A,A-B)\), and

\begin{equation}
\label{eq:ch20-conic-perturbed-degrees}
 \bigl((\theta+\ddc\varphi)\cdot C_0,\ (\theta+\ddc\varphi)\cdot C_1\bigr)
 =(1+B-A,\ 1+A-B).
\end{equation}

This is exactly \eqref{eq:running-perturbed-measure}, now read as a pair of
degrees rather than as a measure, and it reproduces the semipositivity
criterion \eqref{eq:running-nef-interval}: \(\theta+\ddc\varphi\) is
semipositive if and only if \(|A-B|\leq1\).  The total degree is
\((1+B-A)+(1+A-B)=2=(\{\theta\})\), as
\eqref{eq:ch20-special-fibre-degree} requires; note that the perturbation
\(\ddc\varphi\) has total degree zero, which is
\eqref{eq:ch20-curve-laplacian}, second identity, and is the algebraic form
of the statement that a Laplacian has zero total mass
(\eqref{eq:lap-mass-zero}).

Two numerical checks.  With \(A=0,B=1\) the degrees are \((2,0)\), the
extreme semipositive case, and the associated Chambert--Loir measure is
\(2\delta_{v_0}\), which is \eqref{eq:running-perturbed-measure} again.
With \(A=0,B=2\) the degrees are \((3,-1)\): the form is not semipositive,
and the failure is visible on the single vertical curve \(C_1\).  Finally,
with \(A=B\) one recovers \(\ddc\varphi=0\); by
\Cref{thm:ch20-ddc-exact} this is the only way for a model function on this
model to be \(\ddc\)-closed.
\end{example}

\section{The exact sequence of the model differential}
\label{sec:ch20-exact}

The model differential sits in an exact sequence.  This is the single
structural theorem of the chapter: it says that a closed model form is
nothing more than a numerical class on \(X\) together with a model function,
and that the only ambiguity in the second datum is an additive constant.
Every later manipulation of a reference form---changing it by \(\ddc\psi\),
normalizing a potential by its supremum, comparing two references---is an
application of it.

One algebraic input is
needed in the proof.  We use the following form of the extension argument of
K\"unnemann: after a finite extension of the ground field, a numerically
trivial line bundle becomes, up to a positive multiple, algebraically
trivial; its class can then be written as a correspondence from a
numerically trivial line bundle on a smooth projective curve.  We also use
the projection formula for the resulting Chow-group cap product.  These are
standard algebraic intersection-theoretic facts; their precise application
is reproduced below from \cite[proof of Theorem~4.3]{BFJ16}.

We isolate the input so that the reader can see exactly where it is used and
what would be needed to avoid it.

\refstepcounter{lemma}\label{lem:ch20-kunnemann-input}
\begin{quote}
\textbf{Foundational input \thelemma\ (K\"unnemann's correspondence;
source boundary).}
Let \(X\) be a smooth connected projective variety over a field \(K\) and
let \(L\in\Pic(X)\) be numerically trivial.  Then there exist a finite field
extension \(K'/K\), an integer \(m\geq1\), a smooth projective connected
\(K'\)-curve \(T\), a numerically trivial \(\mathbf Q\)-line bundle \(M\) on
\(T\), and a Cartier divisor \(Z\) on \(X_{K'}\times_{K'}T\), such that in
\(\Pic(X_{K'})_{\mathbf Q}\) one has
\(mL_{K'}=q_*(p^*M\cdot Z)\), where \(p,q\) are the two projections.
This precise correspondence is imported from the argument of
\cite[proof of Theorem~4.3]{BFJ16}; no proof is claimed here.  Its
algebraic content rests on the fact that numerical and
algebraic equivalence agree up to torsion for divisors, and on the
representability of the identity component of the Picard scheme by an
abelian variety, which is then generated by images of Jacobians of curves.
\end{quote}
Nothing in \Cref{lem:ch20-kunnemann-input} refers to a valuation, to a
model, or to the residue characteristic.  The reader who is willing to
assume \(N^1(X)=\mathbf R\), as happens for \(X=\mathbf P^n\), may skip
Step~4 of the proof below entirely.

\begin{theorem}\label{thm:ch20-ddc-exact}
Under the standing BFJ hypotheses, the sequence

\begin{equation}
\label{eq:ch20-ddc-exact}
 0\longrightarrow\mathbf R
 \longrightarrow\mathcal D(X)_{\mathbf R}
 \xrightarrow{\ \ddc\ }\mathcal Z^{1,1}(X)
 \xrightarrow{\ \{\cdot\}\ }N^1(X)
 \longrightarrow0
\end{equation}

is exact.
\end{theorem}

\begin{proof}
\emph{Step 0: reduction to a fixed regular model.}
It is enough to prove the following stronger sequence on each regular model:

\begin{equation}
\label{eq:ch20-fixed-model-exact}
 0\longrightarrow \mathbf R\mathcal X_s
 \longrightarrow \Div_0(\mathcal X)_{\mathbf R}
 \longrightarrow N^1(\mathcal X/R)
 \longrightarrow N^1(X)
 \longrightarrow0.
\end{equation}

Indeed, regular models are cofinal under the standing characteristic-zero
hypotheses, and taking the direct limit of
\eqref{eq:ch20-fixed-model-exact} identifies
\(\mathbf R\mathcal X_s\) with the constant model functions by
\eqref{eq:ch20-sign}.

Two remarks on this reduction.  First, cofinality of regular models is
\Cref{prop:snc-cofinal}\,(1), which is where resolution of singularities,
hence \(\operatorname{char}k=0\), is used; \Cref{sec:ch20-scope} records
what is known without it.  Second, a direct limit of exact sequences of
real vector spaces over a directed index category is exact, because
filtered colimits are exact in the category of modules; so no further
argument is needed to pass from \eqref{eq:ch20-fixed-model-exact} to
\eqref{eq:ch20-ddc-exact}.  The transition maps in the second and third
terms are the pullbacks of divisors and classes, which are compatible by
\Cref{lem:ch20-ddc-well-defined}, and the transition maps in the first and
fourth terms are the identity.

\emph{Step 1: exactness at the vertical divisors.}
Write

\begin{equation}
\label{eq:ch20-special-fibre-components}
                       \mathcal X_s=\sum_{i\in I}b_iE_i.
\end{equation}

The special fibre is connected.  In fact, otherwise the theorem on formal
functions would split \(H^0(\mathcal X,\mathcal O_{\mathcal X})\) as a
nontrivial product, whereas its restriction to the generic fibre is a domain
because \(X\) is connected and smooth.  Since \(\mathcal X\) is regular,
each \(E_i\) is Cartier.  Choose a relatively ample line bundle
\(\mathcal A\).  For a vertical divisor \(D=\sum_i a_iE_i\), put

\begin{equation}
\label{eq:ch20-hodge-quadratic}
 q(D):=-\bigl(D^2\cdot\mathcal A^{n-1}\bigr).
\end{equation}

For \(i\ne j\), set

\begin{equation}
\label{eq:ch20-edge-weights}
 w_{ij}:=b_ib_j(E_i\cdot E_j\cdot\mathcal A^{n-1})\geq0.
\end{equation}

The inequality holds because distinct components meet in an effective
cycle and \(\mathcal A\) is relatively ample.  Moreover
\((\mathcal X_s\cdot E_i\cdot\mathcal A^{n-1})=0\), since
\(\mathcal O_{\mathcal X}(\mathcal X_s)\) is trivial.  Writing
\(a_i=b_iu_i\) and using this row-sum identity gives the exact calculation

\begin{equation}
\label{eq:ch20-hodge-laplacian}
 \begin{aligned}
 q(D)
  &=-\sum_{i,j}b_ib_ju_iu_j
       (E_i\cdot E_j\cdot\mathcal A^{n-1})\\
  &=\sum_{i<j}w_{ij}(u_i-u_j)^2.
 \end{aligned}
\end{equation}

Thus \(q\geq0\).  Its nullspace consists of the vectors for which
\(u_i=u_j\) whenever \(E_i\cap E_j\ne\varnothing\).  The incidence graph is
connected because \(\mathcal X_s\) is connected, so all \(u_i\) are equal;
hence \(D\in\mathbf R\mathcal X_s\).  If the numerical class of \(D\)
vanishes, then \(q(D)=0\), proving exactness at
\(\Div_0(\mathcal X)_{\mathbf R}\).

The middle equality in \eqref{eq:ch20-hodge-laplacian} deserves to be
written out once, since the same manipulation recurs in
\Cref{prop:ch20-relative-hodge}, in
\Cref{chap:bfj-compactness}, and in
\Cref{chap:bfj-energy}.  Set \(c_{ij}=(E_i\cdot E_j\cdot\mathcal A^{n-1})\),
so that \(w_{ij}=b_ib_jc_{ij}\) for \(i\ne j\) and
\(\sum_jb_jc_{ij}=0\) for every \(i\).  Then

\[
 \sum_{i,j}b_ib_ju_iu_jc_{ij}
 =\sum_{i,j}b_ib_jc_{ij}\Bigl(
   -\tfrac12(u_i-u_j)^2+\tfrac12u_i^2+\tfrac12u_j^2\Bigr)
 =-\tfrac12\sum_{i\ne j}w_{ij}(u_i-u_j)^2 ,
\]

because \(\sum_{i,j}b_ib_jc_{ij}u_i^2=\sum_ib_iu_i^2\sum_jb_jc_{ij}=0\) and
similarly for the third term.  Multiplying by \(-1\) and using
\(w_{ij}=w_{ji}\) gives \eqref{eq:ch20-hodge-laplacian}.  Note that the
diagonal terms \(c_{ii}\), which are generally negative, never had to be
evaluated: they were eliminated by the row-sum relation.  This is the
mechanism behind every relative Hodge-index computation in the book.

\emph{Step 2: injectivity of \(\mathbf R\mathcal X_s\to\Div_0\).}
This is immediate: \(\mathcal X_s\neq0\) as a divisor, and
\(\varphi_{c\mathcal X_s}=c\) by \eqref{eq:ch20-sign}, so
\(c\mathcal X_s=0\) forces \(c=0\).  Exactness at
\(\mathbf R\mathcal X_s\) and the identification of that line with the
constants are therefore both contained in \eqref{eq:ch20-sign}.

\emph{Step 3: exactness at \(N^1(\mathcal X/R)\), the curve case.}
We next prove exactness at \(N^1(\mathcal X/R)\).  It suffices first to work
over \(\mathbf Q\) and then tensor with \(\mathbf R\).  Let \(L\) be a
numerically trivial line bundle on \(X\).  We must construct a
\(\mathbf Q\)-line bundle \(\mathcal L\) on \(\mathcal X\) that extends
\(L\) and has degree zero on every vertical curve.

Suppose first that \(n=1\).  Choose any extension \(\mathcal L_0\) of
\(L\).  Flatness and the projection formula give

\begin{equation}
\label{eq:ch20-curve-degree-compatibility}
          \sum_i b_i(\mathcal L_0\cdot E_i)=\deg_X L=0.
\end{equation}

By \eqref{eq:ch20-hodge-laplacian}, the symmetric intersection matrix
\(((E_i\cdot E_j))\) has kernel \(\mathbf R(b_i)_i\), hence image the
orthogonal complement of \((b_i)_i\).  The right-hand side
\(((\mathcal L_0\cdot E_i))_i\) belongs to that complement by
\eqref{eq:ch20-curve-degree-compatibility}.  The matrix and the right-hand
side are rational, so there are \(a_i\in\mathbf Q\) such that

\begin{equation}
\label{eq:ch20-curve-vertical-correction}
          \Bigl(\mathcal L_0-\sum_i a_iE_i\Bigr)\cdot E_j=0
          \quad\text{for every }j.
\end{equation}

This proves the curve case.

Two details of this argument are worth making explicit, because the same
two recur in \Cref{chap:bfj-compactness}.  The kernel statement is exactly
Step~1 applied with \(n=1\), \(\mathcal A^{n-1}=1\): the quadratic form
\(q\) of \eqref{eq:ch20-hodge-quadratic} is then
\(-(D^2)\), and \eqref{eq:ch20-hodge-laplacian} identifies it with the
Dirichlet form of the dual graph, whose kernel is the constants.  Since the
matrix is symmetric, its image is the orthogonal complement of its kernel,
which is the hyperplane \(\{\sum_ib_i\lambda_i=0\}\); this is where
\eqref{eq:ch20-curve-degree-compatibility} is used.  Rationality of the
solution \((a_i)\) is automatic because a rational linear system with a
solution over \(\mathbf R\) has one over \(\mathbf Q\).

\emph{Step 4: exactness at \(N^1(\mathcal X/R)\), the general case.}
For arbitrary \(n\), invoke the correspondence input stated before the
theorem, that is, \Cref{lem:ch20-kunnemann-input}.  After replacing \(L\) by
a positive multiple and making a finite
extension \(K'/K\), there are a smooth projective \(K'\)-curve \(T\), a
numerically trivial \(\mathbf Q\)-line bundle \(M\) on \(T\), and a Cartier
divisor \(Z\) on \(X_{K'}\times_{K'}T\) such that

\begin{equation}
\label{eq:ch20-kunnemann-correspondence}
                  L=q_*\bigl(p^*M\cdot Z\bigr)
                  \quad\text{in }\Pic(X)_{\mathbf Q},
\end{equation}

where \(p\) and \(q\) are the two projections.  Let \(R'\) be the integral
closure of \(R\) in \(K'\), choose a regular model \(\mathcal T\) of
\(T\) over \(R'\), and put
\(\mathcal X_{\mathcal T}:=\mathcal X\times_R\mathcal T\).  By the curve
case, \(M\) has a \(\mathbf Q\)-line-bundle extension \(\mathcal M\) on
\(\mathcal T\) that is numerically trivial on its special fibre.  If
\(\mathcal Z\) denotes the closure of \(Z\), define

\begin{equation}
\label{eq:ch20-kunnemann-extension}
             \mathcal L:=q_*\bigl(p^*\mathcal M\cdot\mathcal Z\bigr)
             \in\operatorname{CH}^1(\mathcal X)_{\mathbf Q}
             =\Pic(\mathcal X)_{\mathbf Q}.
\end{equation}

The equality on the right uses regularity of \(\mathcal X\), and restriction
of \eqref{eq:ch20-kunnemann-extension} to the generic fibre recovers
\eqref{eq:ch20-kunnemann-correspondence}.  If \(C\) is any proper vertical
curve on \(\mathcal X\), the Chow projection formula and the module property
of cap products yield

\begin{equation}
\label{eq:ch20-kunnemann-projection}
 \begin{aligned}
  \deg(C\cdot\mathcal L)
   &=\deg\!\left(C\mathbin{\cdot_q}
          (p^*\mathcal M\cdot\mathcal Z)\right)\\
   &=\deg\!\left(p^*\mathcal M\cdot
          (C\mathbin{\cdot_q}\mathcal Z)\right)=0.
 \end{aligned}
\end{equation}

The last equality holds because \(C\mathbin{\cdot_q}\mathcal Z\) is a
vertical one-cycle on \(\mathcal X_{\mathcal T}\), while
\(p^*\mathcal M\) has degree zero on every such curve: if the curve maps to
a point of \(\mathcal T_s\), the restriction is trivial; if it maps onto a
component of \(\mathcal T_s\), the degree is its mapping degree times the
degree of \(\mathcal M\) on that component.  Thus \(\mathcal L\) is the
required relatively numerically trivial extension.

The passage from \(K'\) back to \(K\) is harmless because relative
numerical triviality may be tested after the finite flat base change
\(R'/R\): a vertical curve on \(\mathcal X\) is dominated by a vertical
curve on \(\mathcal X\times_RR'\), and degrees are multiplied by the
corresponding residue-field degree, which is positive.  Likewise the
replacement of \(L\) by a positive multiple is harmless because we work in
\(\mathbf Q\)-coefficients throughout Steps~3 and~4 and only tensor with
\(\mathbf R\) at the end.

\emph{Step 5: the kernel argument.}
Here is the final kernel argument explicitly.  Let
\([\mathcal N]\in N^1(\mathcal X/R)\) restrict to zero in \(N^1(X)\), and
put \(L=\mathcal N|_X\).  The construction supplies a relatively
numerically trivial \(\mathbf Q\)-line bundle \(\mathcal L\) extending
\(L\).  The \(\mathbf Q\)-line bundle
\(\mathcal N\otimes\mathcal L^{-1}\) is trivial on the generic fibre.
After clearing denominators, the divisor of a trivializing rational section
is vertical, say \(D\).  Therefore
\([\mathcal N]=[\mathcal L]+[D]=[D]\) in \(N^1(\mathcal X/R)\).  This proves
that the kernel of restriction is precisely the image of the vertical
divisors.

\emph{Step 6: surjectivity.}
Finally, every class in \(N^1(X)\) is a real linear combination of Cartier
divisor classes.  On a regular model their closures are Cartier, proving
surjectivity of the last arrow in
\eqref{eq:ch20-fixed-model-exact}.  All terms of that sequence are therefore
exact, and passage to the direct limit proves
\eqref{eq:ch20-ddc-exact}.
\end{proof}

\begin{corollary}
\label{cor:ch20-dimension-count}
For a regular model \(\mathcal X\) with \(\#I\) irreducible components in
its special fibre,

\begin{equation}
\label{eq:ch20-dimension-count}
 \dim_{\mathbf R}N^1(\mathcal X/R)
 =\dim_{\mathbf R}N^1(X)+\#I-1 .
\end{equation}
\end{corollary}

\begin{proof}
For a regular model one has
\[
 \dim\Div_0(\mathcal X)_{\mathbf R}=\#I.
\]
Taking alternating dimensions in the exact sequence
\eqref{eq:ch20-fixed-model-exact} now gives the assertion.
\end{proof}

Formula \eqref{eq:ch20-dimension-count} is a convenient audit tool: it
predicted the two dimensions computed by hand in
\Cref{ex:ch20-p2-blowup-numerics}, and it is used again in
\Cref{ex:ch20-two-model-surface}.  It also shows that the vertical part of
\(N^1(\mathcal X/R)\) grows exactly with the combinatorics of the special
fibre, which is why the dual complex, and not the model, is the right
carrier for a potential.

\begin{corollary}\label{cor:ch20-reference-change}
If \(\psi\) is a model function, then

\begin{equation}
\label{eq:ch20-reference-change}
 \varphi\text{ is \(\theta\)-psh at model level}
 \quad\Longleftrightarrow\quad
 \varphi-\psi\text{ is \((\theta+\ddc\psi)\)-psh at model level}.
\end{equation}
\end{corollary}

\begin{proof}
Both curvature forms in \eqref{eq:ch20-reference-change} equal
\(\theta+\ddc\varphi\).
\end{proof}

\begin{corollary}
\label{cor:ch20-normal-form}
Let \(\theta_0\in\mathcal Z^{1,1}(X)\) and let
\(\theta\in\mathcal Z^{1,1}(X)\) satisfy \(\{\theta\}=\{\theta_0\}\).  Then
there is a model function \(\varphi\), unique up to an additive constant,
with \(\theta=\theta_0+\ddc\varphi\).  Consequently, once a reference form
\(\theta_0\) has been fixed in each numerical class, the assignment
\(\varphi\mapsto\theta_0+\ddc\varphi\) identifies

\[
 \mathcal D(X)_{\mathbf R}/\mathbf R
 \;\xrightarrow{\ \sim\ }\;
 \{\theta\in\mathcal Z^{1,1}(X):\{\theta\}=\{\theta_0\}\}.
\]
\end{corollary}

\begin{proof}
\(\theta-\theta_0\) lies in the kernel of \(\{\cdot\}\), hence in the image
of \(\ddc\) by \Cref{thm:ch20-ddc-exact}; uniqueness up to a constant is
exactness at \(\mathcal D(X)_{\mathbf R}\).
\end{proof}

\begin{warning}{Exactness needs regular models, hence characteristic zero}
The proof of \Cref{thm:ch20-ddc-exact} uses regularity of \(\mathcal X\) in
three places: to know that each \(E_i\) is Cartier, so that
\(\Div_0(\mathcal X)_{\mathbf R}=\bigoplus_i\mathbf R E_i\); to identify
\(\operatorname{CH}^1(\mathcal X)_{\mathbf Q}\) with
\(\Pic(\mathcal X)_{\mathbf Q}\) in
\eqref{eq:ch20-kunnemann-extension}; and to have a symmetric intersection
form in Step~3.  Cofinality of regular models is
\Cref{prop:snc-cofinal}, which rests on resolution of singularities and
therefore on \(\operatorname{char}k=0\).  In positive residue
characteristic, where only alterations are available
(\Cref{app:birational-geometry}), the direct-limit statement
\eqref{eq:ch20-ddc-exact} is not proved here.  What survives without
regularity, for an arbitrary normal model, is the inclusion
\(\ddc(\mathcal D(X)_{\mathbf R})\subseteq\ker\{\cdot\}\) and the injectivity
statement of Step~2, both of which are formal.  See
\Cref{sec:ch20-scope}.
\end{warning}

\section{Positive representatives}
\label{sec:ch20-positive-representative}

We can now prove the statement announced in
\Cref{sec:ch20-semipositive-def}: a closed model form with ample generic
class becomes \(\mathcal X\)-positive after the addition of a model
potential, which may be taken arbitrarily small when the form was already
semipositive.  One preliminary lemma is needed, and it is the mechanism by
which relative ampleness with a prescribed generic class propagates upward
through the inverse system of models.

\begin{lemma}
\label{lem:ch20-ample-propagation}
Let \(\mathcal X\) be a model carrying a relatively ample
\(\mathbf R\)-class \(\mathcal L\), and let \(\mathcal Y\) be a model
dominating \(\mathcal X\).  Then there are a model \(\mathcal Z\) dominating
\(\mathcal Y\), with induced domination \(\pi:\mathcal Z\to\mathcal X\), an
effective vertical \(\mathbf Q\)-Cartier divisor \(\mathcal E\) on
\(\mathcal Z\), and \(\delta_0>0\), such that

\begin{equation}
\label{eq:ch20-ample-propagation}
 \pi^*\mathcal L-\delta[\mathcal E]
 \quad\text{is relatively ample on }\mathcal Z
 \quad\text{for }0<\delta<\delta_0 .
\end{equation}

In particular the models carrying a relatively ample \(\mathbf R\)-class
whose generic numerical class equals \(\{\mathcal L|_X\}\) form a cofinal
family.
\end{lemma}

\begin{proof}
The generic fibre \(X\) is a quasi-compact open subscheme of
\(\mathcal X\)---the preimage of the generic point of \(\Spec R\)---and the
domination \(\mathcal Y\to\mathcal X\) is an isomorphism over it.  By
\Cref{thm:appB-admissible-domination} there is a blowup
\(\mathcal Z_0\to\mathcal X\), with centre disjoint from \(X\), which
factors through \(\mathcal Y\).  Let \(\mathcal Z\) be the normalization of
\(\mathcal Z_0\), a normal projective flat \(R\)-model of \(X\) dominating
\(\mathcal Y\), and write \(\pi:\mathcal Z\to\mathcal X\) and
\(\nu:\mathcal Z\to\mathcal Z_0\).

If \(\mathfrak a\subseteq\cO_{\mathcal X}\) is the centre of the blowup,
then \(\mathfrak a\cO_{\mathcal Z_0}=\cO_{\mathcal Z_0}(-\mathcal E_0)\)
with \(\mathcal E_0\) effective Cartier and
\(\cO_{\mathcal Z_0}(-\mathcal E_0)\) ample relative to
\(\mathcal Z_0\to\mathcal X\); this is
\Cref{def:appB-exceptional-divisor} and
\eqref{eq:appB-exceptional-ideal}.  Since the centre is disjoint from
\(X\), the divisor \(\mathcal E_0\) is supported in the special fibre, that
is, vertical.  Put \(\mathcal E:=\nu^*\mathcal E_0\), again effective and
vertical; then
\(\cO_{\mathcal Z}(-\mathcal E)=\nu^*\cO_{\mathcal Z_0}(-\mathcal E_0)\) is
ample relative to \(\mathcal Z\to\mathcal X\), because \(\nu\) is finite and
the pullback of a relatively ample invertible sheaf along a finite morphism
is relatively ample.

Now \(\cO_{\mathcal Z}(-\mathcal E)\) is ample for
\(\pi:\mathcal Z\to\mathcal X\) and \(\mathcal L\) is ample for
\(\mathcal X\to\Spec R\), so by \cite[4.6.13]{EGAII} the class
\(m\pi^*\mathcal L-[\mathcal E]\) is ample for
\(\mathcal Z\to\Spec R\) for all sufficiently large integers \(m\).  Fix
such an \(m\) and set \(\delta_0=1/m\).  For \(0<\delta<\delta_0\) write

\[
 \pi^*\mathcal L-\delta[\mathcal E]
 =\delta\bigl(m\pi^*\mathcal L-[\mathcal E]\bigr)
  +(1-\delta m)\,\pi^*\mathcal L ,
\]

a sum of a positive multiple of a relatively ample class and a positive
multiple of the relatively nef class \(\pi^*\mathcal L\)
(\Cref{prop:nef-pullback}); by
\Cref{lem:ch20-special-fibre-test}\,\cref{itm:ch20-sft:cone} it is relatively
ample.  Since \(\mathcal E\) is vertical, its generic restriction is zero,
so the generic numerical class is unchanged.
\end{proof}

\begin{proposition}\label{prop:ch20-positive-representative}
Let \(\theta\in\mathcal Z^{1,1}(X)\) and assume that
\(\{\theta\}\) is ample.  There is a cofinal family of models
\(\mathcal X\) such that, for each \(\mathcal X\) in the family, a model
function \(\psi\) determined on \(\mathcal X\) makes
\(\theta+\ddc\psi\) \(\mathcal X\)-positive.  If \(\theta\) is
semipositive, then, for every \(\varepsilon>0\), the function \(\psi\) can
be chosen with

\begin{equation}
\label{eq:ch20-small-positive-representative}
                       -\varepsilon\leq\psi\leq0.
\end{equation}
\end{proposition}

\begin{proof}
\emph{Step 1: one model carries a relatively ample class of generic class
\(\{\theta\}\).}  The ample cone of \(X\) is an open convex cone in the
finite-dimensional space \(N^1(X)\), and the rational classes are dense in
\(N^1(X)\).  Choose therefore rational ample classes
\(\beta_1,\ldots,\beta_r\) forming a basis of \(N^1(X)\) and so close to
\(\{\theta\}\) that \(\{\theta\}\) lies in the interior of the cone they
span; then

\begin{equation}
\label{eq:ch20-ample-decomposition}
 \{\theta\}=\sum_{j=1}^rt_j\beta_j,\qquad t_j>0 .
\end{equation}

For each \(j\) pick \(m_j\geq1\) with \(m_j\beta_j\) the class of a very
ample line bundle, embed \(X\hookrightarrow\mathbf P^{N_j}_K\) accordingly,
and let \(\mathcal X_j\subseteq\mathbf P^{N_j}_R\) be the flat closure of the
image (\Cref{chap:models}).  Then \(\mathcal X_j\) is a model and
\(\mathcal B_j:=m_j^{-1}c_1(\cO(1)|_{\mathcal X_j})\) is a relatively ample
\(\mathbf Q\)-class with \(\{\mathcal B_j|_X\}=\beta_j\).  Let
\(\mathcal Y_0\) be a model dominating \(\mathcal X_1,\ldots,\mathcal X_r\)
(\Cref{prop:common-domination}).  Applying
\Cref{lem:ch20-ample-propagation} to \(\mathcal X_1\), \(\mathcal B_1\) and
\(\mathcal Y_0\) produces a model \(\mathcal W\) dominating
\(\mathcal Y_0\), a vertical \(\mathcal E\) on \(\mathcal W\) and
\(\delta\in(0,\delta_0)\) making
\(\pi_1^*\mathcal B_1-\delta[\mathcal E]\) relatively ample on
\(\mathcal W\), where \(\pi_j:\mathcal W\to\mathcal X_j\) denote the induced
dominations.  Then

\[
\begin{aligned}
 \mathcal L_{\mathcal W}
 &:=\sum_{j=1}^rt_j\pi_j^*\mathcal B_j-t_1\delta[\mathcal E]\\
 &\phantom{:}=t_1\bigl(\pi_1^*\mathcal B_1-\delta[\mathcal E]\bigr)
  +\sum_{j\geq2}t_j\pi_j^*\mathcal B_j
\end{aligned}
\]

is relatively ample on \(\mathcal W\), being a positive multiple of a
relatively ample class plus a relatively nef one
(\Cref{lem:ch20-special-fibre-test}\,\cref{itm:ch20-sft:cone}).  Its generic
numerical class is \(\{\theta\}\) by
\eqref{eq:ch20-ample-decomposition}, since \(\mathcal E\) is vertical.

\emph{Step 2: a global correcting potential.}  The element
\([\mathcal L_{\mathcal W}]-\theta\) of \(\mathcal Z^{1,1}(X)\) has generic
class \(\{\theta\}-\{\theta\}=0\), so by exactness of
\eqref{eq:ch20-ddc-exact} at \(\mathcal Z^{1,1}(X)\) there is a model
function \(\psi_0\), unique up to an additive constant, with

\begin{equation}
\label{eq:ch20-positive-difference}
             [\mathcal L_{\mathcal W}]=\theta+\ddc\psi_0 .
\end{equation}

Note that this uses the exact sequence only in the direct limit; no
statement on a fixed model is required.

\emph{Step 3: the cofinal family.}  Let \(\mathcal Y\) be an arbitrary
model.  By \Cref{prop:common-domination} choose \(\mathcal X'\) dominating
\(\mathcal Y\), \(\mathcal W\), a determination of \(\theta\) and a
determination of \(\psi_0\).  By \Cref{lem:ch20-ample-propagation}, applied
to \(\mathcal W\), \(\mathcal L_{\mathcal W}\) and \(\mathcal X'\), there
are a model \(\mathcal Z\) dominating \(\mathcal X'\), with
\(\pi:\mathcal Z\to\mathcal W\), a vertical effective
\(\mathbf Q\)-Cartier divisor \(\mathcal E\) on \(\mathcal Z\), and
\(\delta>0\) such that
\(\pi^*\mathcal L_{\mathcal W}-\delta[\mathcal E]\) is relatively ample on
\(\mathcal Z\).  On \(\mathcal Z\), the classes
\(\pi^*\mathcal L_{\mathcal W}\) and \(\theta_{\mathcal Z}+\ddc\psi_0\) both
determine the same closed model form by
\eqref{eq:ch20-positive-difference}, hence are equal in
\(N^1(\mathcal Z/R)\) by \Cref{cor:ch20-limit-structure}.  Therefore, setting

\[
 \psi:=\psi_0-\delta\varphi_{\mathcal E},
\]

a model function determined on \(\mathcal Z\), the form \(\theta+\ddc\psi\)
has determination \(\pi^*\mathcal L_{\mathcal W}-\delta[\mathcal E]\) on
\(\mathcal Z\) and is therefore \(\mathcal Z\)-positive.  Since
\(\mathcal Z\) dominates the arbitrary \(\mathcal Y\), the models on which
the conclusion holds are cofinal.  By \Cref{prop:snc-cofinal} the collection
may be intersected with the regular models when that is convenient, at the
cost of applying
\Cref{lem:ch20-ample-propagation} once more.  This proves the first
assertion.

\emph{Step 4: smallness when \(\theta\) is semipositive.}  Write
\(\mathcal X\) for the model produced by Step~3 and \(\theta_{\mathcal X}\)
for the determination of \(\theta\) on it.  Adding a constant
to \(\psi\) does not change \(\ddc\psi\), by
\eqref{eq:ch20-ddc-constant}; so we may translate \(\psi\) and assume
\(\max\psi=0\), that is \(\psi\leq0\).  For \(0<t<1\),

\begin{equation}
\label{eq:ch20-positive-convex-combination}
 \theta+\ddc(t\psi)
 =t(\theta+\ddc\psi)+(1-t)\theta ,
\end{equation}

which is relatively ample on \(\mathcal X\): it is the sum of the relatively
ample class \(t(\theta+\ddc\psi)\) and the relatively nef class
\((1-t)\theta_{\mathcal X}\), and
\Cref{lem:ch20-special-fibre-test}\,\cref{itm:ch20-sft:cone} says that such a
sum is relatively ample.  Here \(\theta_{\mathcal X}\) is relatively nef by
\Cref{prop:ch20-semipositive-well-defined}.
Since \(\psi\) is bounded and \(t\psi\leq0\), there is nothing to prove if
\(\psi=0\); otherwise choosing
\(t\leq\varepsilon/\|\psi\|_\infty\) gives
\eqref{eq:ch20-small-positive-representative} with \(t\psi\) in place of
\(\psi\).
\end{proof}

\begin{remark}[What the smallness clause buys]
\label{rem:ch20-smallness}
The second half of \Cref{prop:ch20-positive-representative} is used in
\Cref{chap:bfj-psh,chap:bfj-envelopes} to replace a semipositive reference
form by an \(\mathcal X\)-positive one at the cost of an arbitrarily small
uniform shift of every potential.  Without the uniform bound
\eqref{eq:ch20-small-positive-representative} the substitution would move
sup-norms and energies by an uncontrolled amount, and the compactness
arguments of \Cref{chap:bfj-compactness} would not survive it.  Note also
that \eqref{eq:ch20-positive-convex-combination} is an identity of closed
model forms, not of divisors: the two sides are computed on
\(\mathcal X\), and it is only there that the convex combination makes
sense.  Finally, Step~4 uses semipositivity of \(\theta\) twice---once to
know that \(\theta_{\mathcal X}\) is relatively nef, and once implicitly, in
that \(\psi\leq0\) may be arranged by translation only because
\(\psi\) is bounded, which is automatic for a model function by
\Cref{prop:ch20-maximum-principle}.
\end{remark}

\begin{warning}{The hypothesis is on \(\{\theta\}\), not on \(\theta\)}
\Cref{prop:ch20-positive-representative} assumes that the \emph{generic}
class \(\{\theta\}\) is ample; it assumes nothing about \(\theta\) itself.
This is exactly the right hypothesis for \Cref{part:bfj}, where the
polarization is fixed once and for all and the reference form is allowed to
move within its generic class.  It is also the reason why the statement can
fail to hold with \(\varepsilon\)-smallness when \(\theta\) is not
semipositive: the potential needed to reach the ample cone may then be
large, as \Cref{ex:ch20-nef-generic-not-semipositive} shows, where a shift
of size \(2\) is required.  Nothing in this chapter provides a bound on
\(\|\psi\|_\infty\) in terms of \(\theta\) alone.
\end{warning}

\section{Vertical ideals, maxima, and model envelopes}
\label{sec:ch20-max-envelopes}

\subsection{Vertical ideals}
\label{sec:ch20-vertical-ideals}

The maximum operation is easiest to understand through vertical fractional
ideals.  If \(\mathfrak a\) is such an ideal on \(\mathcal X\), let
\(\rho:\mathcal Y\to\mathcal X\) be the normalized blowup and write

\begin{equation}
\label{eq:ch20-ideal-divisor}
             \mathfrak a\mathcal O_{\mathcal Y}
                    =\mathcal O_{\mathcal Y}(D).
\end{equation}

Then \(\log|\mathfrak a|=\varphi_D\).  If \(\mathfrak a\subset\mathcal O\),
the divisor \(D\) is anti-effective and the function is nonpositive; this is
consistent with \eqref{eq:ch20-sign}.

We recall the conventions, since the signs are the ones that decide whether
a maximum or a minimum appears.  For a coherent fractional ideal
\(\mathfrak a\) on \(\mathcal X\) whose zero locus is contained in
\(\mathcal X_s\), the function \(\log|\mathfrak a|\) on \(X^{\An}\) is
defined by

\begin{equation}
\label{eq:ch20-log-ideal}
 \log|\mathfrak a|(x)=\max_{1\le\alpha\le r}\log|f_\alpha(x)|
 =-\min_{1\le\alpha\le r}v_x(f_\alpha)
\end{equation}

for local generators \(f_1,\dots,f_r\); \Cref{lem:vertical-ideal-convex}
proves that this is independent of the generators, that it is convex and
piecewise affine on every cell of an SNC dual complex, and that it equals
\(-\varphi_D\) when \(\mathfrak a=\cO(-D)\) is invertible with \(D\)
vertical effective.  Because the maximum in \eqref{eq:ch20-log-ideal} is
over logarithms, \emph{sums of ideals produce maxima of functions}:

\begin{equation}
\label{eq:ch20-ideal-sum-max}
 \log|\mathfrak a+\mathfrak b|
 =\max\{\log|\mathfrak a|,\log|\mathfrak b|\},
 \qquad
 \log|\mathfrak a\cdot\mathfrak b|
 =\log|\mathfrak a|+\log|\mathfrak b| .
\end{equation}

This is the entire mechanism of \Cref{prop:ch20-model-max}: the maximum of
two model functions is realized by a sum of ideals, and positivity is then
read off from global generation.

\begin{lemma}\label{lem:ch20-generated-ideal-psh}
Let \(\mathcal L\) be a line bundle on a model \(\mathcal X\), and let
\(\theta=c_1(\mathcal L,\|\cdot\|_{\mathcal L})\).  If
\(\mathcal L\otimes\mathfrak a\) is relatively globally generated, then
\(\log|\mathfrak a|\) is a \(\theta\)-psh model function.
\end{lemma}

\begin{proof}
On the normalized blowup in \eqref{eq:ch20-ideal-divisor}, the line bundle
\(\rho^*\mathcal L+D\) is globally generated, hence relatively nef.  Its
class is \(\theta+\ddc\log|\mathfrak a|\).
\end{proof}

The two implications hidden in that one-line proof are worth naming.  A
relatively globally generated line bundle is relatively nef, because a
global generation gives a morphism to a projective space over the base under
which the bundle is a pullback of \(\cO(1)\), and pullbacks of nef classes
are nef.  And the identification of the class of \(\rho^*\mathcal L+D\) with
\(\theta+\ddc\log|\mathfrak a|\) is
\Cref{lem:ch20-ddc-well-defined} together with
\(\log|\mathfrak a|=\varphi_D\) on the normalized blowup.  Neither step uses
positivity of \(\mathfrak a\) itself; it is the \emph{product}
\(\mathcal L\otimes\mathfrak a\) that has to be positive.

\subsection{Stability under maxima}
\label{sec:ch20-max}

\begin{proposition}\label{prop:ch20-model-max}
The real \(\theta\)-psh model functions form a convex set and are stable
under finite maxima.
\end{proposition}

\begin{proof}
Convexity is immediate from the convexity of the nef cone.  The maximum
requires more care.  First take rational model functions
\(\varphi_1,\varphi_2\) and choose a common determination
\(\mathcal X\), with vertical divisors \(D_1,D_2\).  Represent
\(\theta_{\mathcal X}\) by

\begin{equation}
\label{eq:ch20-ample-basis}
                    \mathcal L=\sum_{j=1}^r t_j\mathcal A_j,
\end{equation}

where the \(\mathcal A_j\) are ample line bundles whose classes span
\(N^1(\mathcal X/R)\).  Choose small positive \(\varepsilon_j\) so that
the coefficients \(t_j+\varepsilon_j\) are rational.  Since
\(\mathcal L+D_i\) is nef,

\begin{equation}
\label{eq:ch20-ample-perturbation}
              \mathcal L_\varepsilon+D_i,
              \qquad
              \mathcal L_\varepsilon
                :=\sum_j(t_j+\varepsilon_j)\mathcal A_j,
\end{equation}

is ample for \(i=1,2\).  After multiplying by \(m\), both corresponding
line bundles are globally generated.  Set

\begin{equation}
\label{eq:ch20-max-ideal}
 \mathfrak a_m
   :=\mathcal O_{\mathcal X}(mD_1)
       +\mathcal O_{\mathcal X}(mD_2).
\end{equation}

Then

\begin{equation}
\label{eq:ch20-max-ideal-function}
             \log|\mathfrak a_m|
                  =m\max\{\varphi_1,\varphi_2\},
\end{equation}

and \(m\mathcal L_\varepsilon\otimes\mathfrak a_m\) is globally
generated.  By \Cref{lem:ch20-generated-ideal-psh}, the maximum is psh with
respect to
\(\theta+\sum_j\varepsilon_jc_1(\mathcal A_j)\).  Letting all
\(\varepsilon_j\) tend to zero and using closedness of the nef cone proves
the rational case.  Real model functions follow by rational approximation
on a common finite-dimensional space and another use of closedness.
\end{proof}

For later use we spell out the four steps that were compressed there, since
each one is a place where the argument could be misremembered.

\begin{enumerate}
\item \emph{The perturbation is needed because \(\theta_{\mathcal X}\) is a
 real class.}  Global generation is a statement about honest line bundles.
 Writing \(\theta_{\mathcal X}\) in the basis \eqref{eq:ch20-ample-basis}
 and rounding the coefficients up to rationals produces
 \(\mathcal L_\varepsilon\) with \(\mathcal L_\varepsilon-\mathcal L\)
 ample, so that \(\mathcal L_\varepsilon+D_i\), being the sum of the nef
 class \(\mathcal L+D_i\) and an ample class, is ample by
 \Cref{lem:ch20-special-fibre-test}\,\cref{itm:ch20-sft:cone}.
\item \emph{Global generation of a large multiple.}  On a projective
 \(R\)-scheme, a relatively ample \(\mathbf Q\)-line bundle has a multiple
 that is relatively very ample, hence relatively globally generated.  We
 choose one \(m\) that works simultaneously for \(i=1,2\) and clears the
 denominators of \(D_1,D_2\) and of the \(t_j+\varepsilon_j\).
\item \emph{The ideal computes the maximum.}  Identity
 \eqref{eq:ch20-max-ideal-function} is the first half of
 \eqref{eq:ch20-ideal-sum-max} together with
 \(\log|\cO(mD_i)|=m\varphi_{D_i}\), which is
 \eqref{eq:ch20-sign} and additivity.  Note that
 \(\mathfrak a_m\) is a \emph{fractional} ideal unless \(D_1,D_2\leq0\);
 nothing in \Cref{lem:ch20-generated-ideal-psh} requires an honest ideal.
\item \emph{Global generation of the twisted ideal.}  The sheaf
 \(m\mathcal L_\varepsilon\otimes\mathfrak a_m\) is generated by the
 sections of the two globally generated summands
 \(m\mathcal L_\varepsilon+mD_i\), because a sum of ideals is generated by
 the union of their generators.  This is the only place where the sum, and
 not the product, of the two ideals is used, and it is the reason maxima
 rather than sums of potentials come out.
\end{enumerate}

\begin{remark}[Why the naive argument fails]
\label{rem:ch20-why-max-is-hard}
This proof is worth remembering.  The maximum is not obtained by raising
one coefficient of a nef vertical divisor, an operation that can leave the
nef cone.  Global generation after an ample perturbation is the mechanism
that preserves positivity.

The precise obstruction is visible in the vertex bookkeeping.  Fix a regular
model with \(\mathcal X_s=\sum_ib_iE_i\), let \(\theta\) be determined on
it, and describe a vertical divisor by its vertex values
\(u_i=\varphi_D(x_{E_i})\).  For a vertical curve \(C\), set
\(\lambda_i=b_i(E_i\cdot C)\).  Then \(\sum_i\lambda_i=0\) because
\(\cO(\mathcal X_s)\simeq\cO\), and

\begin{equation}
\label{eq:ch20-curve-lambda}
 (D\cdot C)=\sum_i\lambda_iu_i
 =\sum_{i}\lambda_i\,(u_i-u_{i_0})
 \qquad\text{for any index }i_0 .
\end{equation}

Since \(\lambda_i\geq0\) whenever \(C\not\subseteq E_i\), the vector
\(\lambda\) has at most \(\#\{i:C\subseteq E_i\}\) negative entries.  Suppose
that for a given \(C\) there is at most one index \(i_0\) with
\(\lambda_{i_0}<0\).  Then, taking that index in
\eqref{eq:ch20-curve-lambda} and interchanging \(u\) and \(u'\) if
necessary so that \(m_{i_0}=u_{i_0}\) for the componentwise maximum
\(m=\max\{u,u'\}\), one gets \(m_i-m_{i_0}\geq u_i-u_{i_0}\) for every
\(i\), hence

\[
 (D_m\cdot C)=\sum_{i\ne i_0}\lambda_i(m_i-m_{i_0})
 \geq\sum_{i\ne i_0}\lambda_i(u_i-u_{i_0})=(D_u\cdot C)
 \geq-(\theta\cdot C),
\]

so the constraint attached to \(C\) is inherited by \(m\).  If instead two or
more \(\lambda_i\) are negative---which requires \(C\) to lie in two
components on each of which it has negative self-intersection---the
monotonicity argument breaks down, and no elementary substitute is known to
us.  This is why the proof of \Cref{prop:ch20-model-max} goes through global
generation.  On a curve every vertical curve is a whole component, so the
favourable case always holds and the argument does succeed
(\Cref{prop:ch20-curve-lattice}).  Even then a second obstruction remains,
independent of positivity: in dimension \(n\ge2\) the maximum of two model
functions determined on \(\mathcal X\) is usually not determined on
\(\mathcal X\) (\Cref{ex:ch20-max-fails-naively}).
\end{remark}

\begin{proposition}
\label{prop:ch20-curve-lattice}
Let \(n=1\), let \(\mathcal X\) be a regular model of \(X\) with
\(\mathcal X_s=\sum_{i\in I}b_iE_i\), and let \(\theta\) be determined on
\(\mathcal X\), with \(d_i=(\theta\cdot E_i)\).  Put

\begin{equation}
\label{eq:ch20-curve-admissible-set}
 C_\theta(\mathcal X):=
 \Bigl\{u\in\mathbf R^I:\ d_i+\sum_{j\neq i}b_j(E_j\cdot E_i)(u_j-u_i)
   \geq0\ \text{ for all }i\in I\Bigr\}.
\end{equation}

Then \(C_\theta(\mathcal X)\) is a closed convex subset of \(\mathbf R^I\),
invariant under adding constants, and stable under componentwise maximum.  A
model function determined on \(\mathcal X\) is \(\theta\)-psh if and only if
its vertex vector lies in \(C_\theta(\mathcal X)\).  It is \emph{not} stable
under componentwise minimum in general
\textup{(\Cref{ex:ch20-min-fails})}.
\end{proposition}

\begin{proof}
The description of \(\theta\)-pshness is
\Cref{prop:ch20-curve-laplacian}: the only vertical curves are the
components \(E_i\) themselves, and \(((\theta+\ddc\varphi_D)\cdot E_i)\) is
the left-hand side of \eqref{eq:ch20-curve-admissible-set}.  Closedness,
convexity and invariance under constants are clear, the last because the
increments \(u_j-u_i\) are unchanged.

Stability under maxima.  Let \(u,u'\in C_\theta(\mathcal X)\), put
\(m=\max\{u,u'\}\) componentwise, and fix \(i\).  Interchanging the roles of
\(u\) and \(u'\) if necessary, assume \(m_i=u_i\).  For every \(j\) we have
\(m_j\geq u_j\), hence \(m_j-m_i\geq u_j-u_i\).  Since the weights
\(b_j(E_j\cdot E_i)\), \(j\neq i\), are nonnegative
(\Cref{prop:ch20-curve-laplacian}),

\begin{equation}
\label{eq:ch20-curve-max-monotone}
 d_i+\sum_{j\neq i}b_j(E_j\cdot E_i)(m_j-m_i)
 \geq
 d_i+\sum_{j\neq i}b_j(E_j\cdot E_i)(u_j-u_i)\geq0 .
\end{equation}

As \(i\) was arbitrary, \(m\in C_\theta(\mathcal X)\).
\end{proof}

The asymmetry between maxima and minima in
\Cref{prop:ch20-curve-lattice} is not an artefact of the proof.  It is the
reason why the pluripotential theory of
\Crefrange{chap:bfj-psh}{chap:bfj-envelopes} is built on suprema and upper
semicontinuous regularization, and why decreasing limits are handled by a
separate mechanism.

\begin{example}[The componentwise minimum leaves the admissible set]
\label{ex:ch20-min-fails}
Let \(n=1\) and let \(\mathcal X\) be a regular semistable model whose dual
graph is a path with three vertices \(E_1-E_2-E_3\), all multiplicities one
and all adjacent intersection numbers one; thus \((E_1^2)=(E_3^2)=-1\) and
\((E_2^2)=-2\).  Take \(\theta\) with degrees \((d_1,d_2,d_3)=(1,0,1)\), so
that \(\{\theta\}\) has degree \(2\).  By
\eqref{eq:ch20-curve-admissible-set} the admissible vertex vectors are

\[
 C_\theta(\mathcal X)=
 \bigl\{u\in\mathbf R^3:\ 1+u_2-u_1\geq0,\quad
   u_1+u_3-2u_2\geq0,\quad 1+u_2-u_3\geq0\bigr\}.
\]

The vectors \(u=(0,0,0)\) and \(u'=(1,0,-1)\) both lie in
\(C_\theta(\mathcal X)\): for \(u\) the three left-hand sides are
\((1,0,1)\), and for \(u'\) they are \((0,0,2)\).  Their componentwise
maximum \((1,0,0)\) lies in \(C_\theta(\mathcal X)\), with left-hand sides
\((0,1,1)\), as \Cref{prop:ch20-curve-lattice} guarantees.  Their
componentwise minimum is \(m'=(0,0,-1)\), for which the middle left-hand
side is \(0+(-1)-0=-1<0\).  Hence \(m'\notin C_\theta(\mathcal X)\), and the
model function with vertex values \(m'\) is not \(\theta\)-psh.

Geometrically: \(u\) is the constant function and \(u'\) is affine of slope
\(-1\); both are ``convex enough'' at the middle vertex, one with equality.
Their minimum has a strictly concave corner there, and a concave corner is
exactly what \eqref{eq:ch20-curve-admissible-set} forbids.
\end{example}

\begin{example}[The vertex-value description on a blown-up plane, and where
the naive arguments run out]
\label{ex:ch20-max-fails-naively}
Return to \(\mathcal X_1=\operatorname{Bl}_p\mathbf P^2_R\) of
\Cref{ex:ch20-p2-blowup-numerics}, with components \(F_1,E\).  Here are the
vectors \(\lambda=(\lambda_{F_1},\lambda_E)\) of
\eqref{eq:ch20-curve-lambda} for the four kinds of vertical curve; all
multiplicities are one, so \(\lambda_i=(E_i\cdot C)\).  Recall
\(\cO(E)|_E=\cO_{\mathbf P^2}(-1)\) and
\(\cO(F_1)|_{F_1}=\cO_{F_1}(-e)\), hence also
\(\cO(F_1)|_E=\cO(\rho^*F-E)|_E=\cO_{\mathbf P^2}(1)\).

\begin{center}
\small
\begin{tabular}{p{.34\textwidth}|c|c|c}
\toprule
\textbf{vertical curve \(C\)}&\(\lambda_{F_1}\)&\(\lambda_E\)
 &\((D_u\cdot C)\)\\
\midrule
\(e=F_1\cap E\), a line of \(E\)&\(+1\)&\(-1\)&\(u_{F_1}-u_E\)\\
\(\ell_E\subseteq E\) any other line&\(+1\)&\(-1\)&\(u_{F_1}-u_E\)\\
\(\tilde\ell\subseteq F_1\) of class \(h_1\)&\(0\)&\(0\)&\(0\)\\
\(\tilde m\subseteq F_1\) of class \(h_1-e\)&\(-1\)&\(+1\)&\(u_E-u_{F_1}\)\\
\bottomrule
\end{tabular}
\end{center}

The last row is the strict transform of a line through \(p\):
\((E\cdot\tilde m)=(e\cdot(h_1-e))_{F_1}=0-(-1)=1\) and
\((F_1\cdot\tilde m)=(-e\cdot(h_1-e))_{F_1}=-1\).  Every row sums to zero,
as it must.

Two consequences.  First, each \(\lambda\) has exactly one negative entry,
so by \Cref{rem:ch20-why-max-is-hard} the componentwise maximum of two
admissible vertex vectors is again admissible on this model.  This is not an
accident of the example: with only two components, the vertex vectors modulo
the translation line form a one-dimensional space parameterized by
\(\delta=u_{F_1}-u_E\), the admissible set is an interval
\(-(\theta\cdot e)\le\delta\le(\theta\cdot\tilde m)\), and componentwise
maximum moves \(\delta\) into the interval spanned by \(\delta,\delta'\).
The genuine obstruction to a monotonicity argument requires a vertical curve
lying in two components with negative self-intersection in each, which
cannot happen when \(\#I=2\) and \(E\cong\mathbf P^2\).

Second---and this is the obstruction that always occurs---the maximum
\(\max\{\varphi_1,\varphi_2\}\) of two model functions determined on
\(\mathcal X\) is, in dimension \(n\geq2\), in general \emph{not} determined
on \(\mathcal X\), whatever its vertex values.  Its restriction to a face of
\(\Delta_{\mathcal X}\) of positive dimension is a maximum of two affine
functions, hence convex and piecewise affine but not affine, so by
\Cref{prop:vertical-divisor-affine-iso} it is not the model function of any
divisor on \(\mathcal X\); and it is strictly larger than the affine
interpolation of its own vertex values.  The blowup implicit in
\eqref{eq:ch20-ideal-divisor} is precisely the subdivision that introduces
the missing breakpoints, and \Cref{ex:ch20-nodal-max} exhibits one on the
nodal interval.  A statement about vertex vectors on a fixed model, such as
\Cref{prop:ch20-curve-lattice}, therefore does not by itself prove
\Cref{prop:ch20-model-max}, even in the favourable sign situation.
\end{example}

\begin{corollary}
\label{cor:ch20-truncation}
Let \(\theta\) be semipositive.  Then every constant is a \(\theta\)-psh
model function, and for every \(\theta\)-psh model function \(\varphi\) and
every \(c\in\mathbf R\) the truncation \(\max\{\varphi,c\}\) is again a
\(\theta\)-psh model function.
\end{corollary}

\begin{proof}
Constants: \(\ddc c=0\) by \eqref{eq:ch20-ddc-constant}, so
\(\theta+\ddc c=\theta\) is semipositive.  Truncation is then
\Cref{prop:ch20-model-max}.
\end{proof}

\begin{proposition}
\label{prop:ch20-maximum-principle}
Let \(\varphi\) be a model function determined on an SNC model
\(\mathcal X\).  Then

\begin{equation}
\label{eq:ch20-maximum-principle}
 \max_{X^{\An}}\varphi=\max_{i\in I}\varphi(x_{E_i}),
 \qquad
 \min_{X^{\An}}\varphi=\min_{i\in I}\varphi(x_{E_i}),
\end{equation}

both extrema being attained at divisorial points of \(\mathcal X\).  In
particular \(\|\varphi\|_\infty=\max_i|\varphi(x_{E_i})|\).
\end{proposition}

\begin{proof}
By \Cref{lem:model-function-affine}, \(\varphi=\varphi\circ r_{\mathcal X}\),
so the extrema over \(X^{\An}\) equal the extrema over the compact dual
complex \(\Delta_{\mathcal X}\).  On \(\Delta_{\mathcal X}\), \(\varphi\) is
affine on each face; an affine function on a simplex attains its extrema at
vertices, and the vertices of \(\Delta_{\mathcal X}\) are the divisorial
points \(x_{E_i}\).
\end{proof}

\subsection{Model envelopes}
\label{sec:ch20-envelopes}

The maximum construction of \Cref{prop:ch20-model-max} makes the family of
\(\theta\)-psh model functions below a given obstacle upward directed, and
this is exactly what is needed to define an envelope.  We record here the
properties that can be established at model level, and we are explicit about
those that cannot.

\begin{definition}[Model envelope]
\index{envelope!model}
\label{def:ch20-model-envelope}
Let \(\theta\in\mathcal Z^{1,1}(X)\) be semipositive and let
\(\varphi:X^{\An}\to\mathbf R\) be bounded.  The \emph{model envelope} of
\(\varphi\) is

\begin{equation}
\label{eq:ch20-model-envelope}
 P^{\model}_\theta\varphi
 :=\sup\bigl\{\psi\in\mathcal D(X)_{\mathbf R}:\
   \psi\ \theta\text{-psh at model level},\ \psi\leq\varphi\bigr\},
\end{equation}

the supremum being taken pointwise on \(X^{\An}\).
\end{definition}

\begin{theorem}
\label{thm:ch20-model-envelope}
Assume \eqref{eq:ch20-standing} and let \(\theta\) be semipositive.  For
bounded functions \(\varphi,\varphi'\) on \(X^{\An}\) and \(c\in\mathbf R\),
\(t\in[0,1]\):
\begin{enumerate}
\item\label{itm:ch20-env:directed} the family in
 \eqref{eq:ch20-model-envelope} is nonempty and upward directed, and
 \(P^{\model}_\theta\varphi\) is the pointwise limit of an increasing net of
 \(\theta\)-psh model functions;
\item\label{itm:ch20-env:bounds}
 \(\inf_{X^{\An}}\varphi\leq P^{\model}_\theta\varphi\leq\varphi\);
\item\label{itm:ch20-env:monotone} \(\varphi\leq\varphi'\) implies
 \(P^{\model}_\theta\varphi\leq P^{\model}_\theta\varphi'\);
\item\label{itm:ch20-env:translate}
 \(P^{\model}_\theta(\varphi+c)=P^{\model}_\theta\varphi+c\);
\item\label{itm:ch20-env:lipschitz}
 \(\bigl\|P^{\model}_\theta\varphi-P^{\model}_\theta\varphi'\bigr\|_\infty
   \leq\|\varphi-\varphi'\|_\infty\);
\item\label{itm:ch20-env:idem}
 \(P^{\model}_\theta(P^{\model}_\theta\varphi)=P^{\model}_\theta\varphi\);
 and \(P^{\model}_\theta\varphi=\varphi\) if \(\varphi\) is a
 \(\theta\)-psh model function;
\item\label{itm:ch20-env:concave}
 \(P^{\model}_\theta(t\varphi+(1-t)\varphi')
   \geq tP^{\model}_\theta\varphi+(1-t)P^{\model}_\theta\varphi'\);
\item\label{itm:ch20-env:reference} for a model function \(\psi\) such that
 \(\theta+\ddc\psi\) is semipositive,
 \(P^{\model}_{\theta+\ddc\psi}(\varphi-\psi)
    =P^{\model}_\theta\varphi-\psi\).
\end{enumerate}
\end{theorem}

\begin{proof}
\cref{itm:ch20-env:directed}.  Nonemptiness: since \(\theta\) is
semipositive, the constant \(\inf\varphi\) is a \(\theta\)-psh model
function \(\leq\varphi\), by \Cref{cor:ch20-truncation}.  Directedness: if
\(\psi_1,\psi_2\) belong to the family, then so does
\(\max\{\psi_1,\psi_2\}\), by \Cref{prop:ch20-model-max}, and it dominates
both.  An upward directed family has a pointwise supremum realized as the
limit of the increasing net indexed by the family itself, ordered by
\(\leq\); nets, and not sequences, are used here because the family is not
countable and no metrizability has been proved
(\Cref{app:nets}).

\cref{itm:ch20-env:bounds}.  The lower bound is the constant just exhibited;
the upper bound is the constraint \(\psi\leq\varphi\).

\cref{itm:ch20-env:monotone}.  The defining family for \(\varphi\) is
contained in that for \(\varphi'\).

\cref{itm:ch20-env:translate}.  \(\psi\mapsto\psi+c\) is a bijection of the
\(\theta\)-psh model functions, by \eqref{eq:ch20-ddc-constant}, and it
matches the constraint \(\psi\leq\varphi\) with \(\psi+c\leq\varphi+c\).

\cref{itm:ch20-env:lipschitz}.  Put \(\delta=\|\varphi-\varphi'\|_\infty\).
Then \(\varphi\leq\varphi'+\delta\), so by
\cref{itm:ch20-env:monotone} and \cref{itm:ch20-env:translate},
\(P^{\model}_\theta\varphi\leq P^{\model}_\theta\varphi'+\delta\);
symmetrically for the other inequality.

\cref{itm:ch20-env:idem}.  Write \(P=P^{\model}_\theta\).  Every \(\psi\) in
the family for \(\varphi\) satisfies \(\psi\leq P\varphi\), so the family
for \(\varphi\) is contained in that for \(P\varphi\); conversely
\(P\varphi\leq\varphi\) gives the reverse inclusion by
\cref{itm:ch20-env:monotone}.  Hence \(P(P\varphi)=P\varphi\).  If
\(\varphi\) is itself a \(\theta\)-psh model function it belongs to its own
family, so \(P\varphi\geq\varphi\), and \cref{itm:ch20-env:bounds} gives
equality.

\cref{itm:ch20-env:concave}.  If \(\psi\leq\varphi\) and
\(\psi'\leq\varphi'\) are \(\theta\)-psh model functions then
\(t\psi+(1-t)\psi'\) is a \(\theta\)-psh model function
(\Cref{prop:ch20-model-max}, convexity) bounded by
\(t\varphi+(1-t)\varphi'\).  Take suprema.

\cref{itm:ch20-env:reference}.  By \Cref{cor:ch20-reference-change},
\(\psi'\mapsto\psi'-\psi\) is a bijection from the \(\theta\)-psh model
functions onto the \((\theta+\ddc\psi)\)-psh model functions, and it carries
the constraint \(\psi'\leq\varphi\) to \(\psi'-\psi\leq\varphi-\psi\).
\end{proof}

\begin{warning}{What the model envelope is not}
\Cref{thm:ch20-model-envelope} asserts nothing about the regularity of
\(P^{\model}_\theta\varphi\).  In general it is neither a model function nor
continuous, and it need not even be upper semicontinuous: it is an
increasing supremum, hence lower semicontinuous when \(\varphi\) is, and
that is the wrong side.  The statement that
\((P^{\model}_\theta\varphi)^*\) is \(\theta\)-psh in the sense of
\Cref{chap:bfj-psh}, and the statement that
\(P^{\model}_\theta\varphi\) is continuous when \(\varphi\) is, are theorems
that require the compactness and regularization machinery of
\Cref{chap:bfj-compactness} and the multiplier-ideal arguments of
\Cref{chap:bfj-envelopes}; they are not consequences of anything in this
chapter.  Nor is the model envelope the same as the psh envelope
\(P_\theta\varphi\) of \Cref{chap:bfj-envelopes}, in which the supremum runs
over all \(\theta\)-psh functions.  The inequality
\(P^{\model}_\theta\varphi\leq P_\theta\varphi\) is immediate; equality is
the content of the decreasing-regularization theorem of
\Cref{chap:bfj-compactness}, again for a restricted class of \(\varphi\).
\end{warning}

\begin{proposition}
\label{prop:ch20-curve-envelope}
Let \(n=1\), let \(\mathcal X\) be a regular model of \(X\), and let
\(\theta\) be determined on \(\mathcal X\) and semipositive, so that
\(d_i=(\theta\cdot E_i)\geq0\).  Let \(g\in\mathbf R^I\) and set

\[
 C_\theta(\mathcal X)_{\le g}
 :=\{u\in C_\theta(\mathcal X):u\leq g\ \text{componentwise}\}
\]

with \(C_\theta(\mathcal X)\) as in
\eqref{eq:ch20-curve-admissible-set}.  Then
\(C_\theta(\mathcal X)_{\le g}\) is nonempty and has a largest element
\(u^\star\), and the corresponding model function
\(\varphi_{u^\star}\) is the largest \(\theta\)-psh model function
determined on \(\mathcal X\) whose vertex values are bounded by \(g\).
\end{proposition}

\begin{proof}
Nonemptiness: the constant vector \((\min_ig_i)\mathbf 1\) lies in
\(C_\theta(\mathcal X)\) because \(d_i\geq0\), and it is \(\leq g\).
Boundedness above is the constraint \(u\leq g\).  By
\Cref{prop:ch20-curve-lattice} the set \(C_\theta(\mathcal X)_{\le g}\) is
closed, convex and stable under componentwise maximum; being a nonempty
subset of \(\mathbf R^I\) with \(I\) finite, bounded above and stable under
finite componentwise maxima, its componentwise supremum
\(u^\star_i:=\sup\{u_i:u\in C_\theta(\mathcal X)_{\le g}\}\) is finite for
each \(i\) and is a limit of an increasing sequence of elements of the set
(take maxima of finitely many near-optimal vectors).  Closedness gives
\(u^\star\in C_\theta(\mathcal X)_{\le g}\), and \(u^\star\) dominates every
element by construction.
\end{proof}

\begin{remark}[Why the fixed-model envelope is a curve phenomenon]
\label{rem:ch20-fixed-model-envelope-scope}
\Cref{prop:ch20-curve-envelope} solves the envelope problem inside a fixed
model, by a finite-dimensional argument, and the solution is again a model
function on the same model.  In dimension \(n\geq2\) neither conclusion is
available by this route: \Cref{ex:ch20-max-fails-naively} shows that the
admissible set of vertex vectors is not stable under componentwise maximum,
and even if it were, \Cref{prop:ch20-model-max} produces the maximum only
after a blowup.  This is precisely why
\Cref{chap:bfj-envelopes} needs multiplier ideals rather than convex
geometry, and why continuity of envelopes is a theorem with hypotheses
rather than a formality.  It is also why the theory does not simply take
place on one dual complex.
\end{remark}

\section{Nefness, semipositivity, and numerical positivity}
\label{sec:ch20-positivity}

The word \emph{positive} has been used in four senses so far.  This section
separates them, proves the implications that hold, and gives a
counterexample or a rigidity theorem for each implication that fails.  The
distinctions are not pedantic: \Cref{chap:bfj-psh} defines
\(\PSH(X,\theta)\) by a condition that refers to semipositivity of
\(\theta+\ddc\varphi\), \Cref{chap:bfj-envelopes} needs
\(\mathcal X\)-positivity of the reference form, and
\Cref{chap:bfj-ma} needs nefness of the generic class to have a positive
total mass.

\begin{definition}[Four positivity notions]
\index{positivity!four model notions}\index{numerical positivity}
\label{def:ch20-four-positivities}
Let \(\theta\in\mathcal Z^{1,1}(X)\) and let \(\mathcal X\) be a model on
which \(\theta\) is determined by \(\theta_{\mathcal X}\).
\begin{enumerate}
\item\label{itm:ch20-pos:ample} \(\theta\) is \emph{\(\mathcal X\)-positive}
 if \(\theta_{\mathcal X}\) is relatively ample.
\item\label{itm:ch20-pos:numerical} \(\theta_{\mathcal X}\) is
 \emph{numerically positive on the special fibre} if
 \((\theta_{\mathcal X}^{\dim Z}\cdot Z)>0\) for every integral closed
 subscheme \(Z\subseteq\mathcal X_s\) with \(\dim Z\geq1\).
\item\label{itm:ch20-pos:semi} \(\theta\) is \emph{semipositive} if
 \(\theta_{\mathcal X}\) is relatively nef
 (\Cref{def:ch20-semipositive}).
\item\label{itm:ch20-pos:generic} \(\theta\) has \emph{nef generic class}
 if \(\{\theta\}\) is nef in \(N^1(X)\).
\end{enumerate}
\end{definition}

\begin{theorem}
\label{thm:ch20-positivity-implications}
With notation as in \Cref{def:ch20-four-positivities}:
\begin{enumerate}
\item\label{itm:ch20-impl:aton}
 \cref{itm:ch20-pos:ample}\(\Rightarrow\)\cref{itm:ch20-pos:numerical}, and
 the converse holds when \(\theta_{\mathcal X}\) is a
 \(\mathbf Q\)-class.
\item\label{itm:ch20-impl:atos}
 \cref{itm:ch20-pos:ample}\(\Rightarrow\)\cref{itm:ch20-pos:semi}.  The
 converse fails: \Cref{ex:ch20-nef-not-ample}.
\item\label{itm:ch20-impl:stog}
 \cref{itm:ch20-pos:semi}\(\Rightarrow\)\cref{itm:ch20-pos:generic}
 \textup{(\Cref{cor:ch20-semipositive-generic-nef})}.  The converse fails
 in the strongest possible way: see
 \Cref{thm:ch20-vertical-nef-rigidity} and
 \Cref{ex:ch20-nef-generic-not-semipositive}.
\item\label{itm:ch20-impl:ntos}
 \cref{itm:ch20-pos:numerical}\(\Rightarrow\)\cref{itm:ch20-pos:semi}
 for a \(\mathbf Q\)-class, by \cref{itm:ch20-impl:aton} and
 \cref{itm:ch20-impl:atos}; but \cref{itm:ch20-pos:numerical} is not stable
 under domination, whereas \cref{itm:ch20-pos:semi} is.
\end{enumerate}
\end{theorem}

\begin{proof}
\cref{itm:ch20-impl:aton}.  If \(\theta_{\mathcal X}\) is relatively ample
then, by
\Cref{lem:ch20-special-fibre-test}\,\cref{itm:ch20-sft:ample}, its
restriction to \(\mathcal X_s\) is ample; the restriction to an integral
closed \(Z\subseteq\mathcal X_s\) is then ample, and the top
self-intersection of an ample class on an integral projective scheme of
dimension \(d\geq1\) over a field is positive.  Conversely, for a
\(\mathbf Q\)-class, numerical positivity on all integral
\(Z\subseteq\mathcal X_s\) is the hypothesis of the Nakai--Moishezon
criterion on the projective \(k\)-scheme \(\mathcal X_s\), which gives
ampleness of the restriction, hence relative ampleness by
\Cref{lem:ch20-special-fibre-test}\,\cref{itm:ch20-sft:ample} again.

\cref{itm:ch20-impl:atos}.  Ample implies nef.  The failure of the converse
is \Cref{ex:ch20-nef-not-ample}.

\cref{itm:ch20-impl:stog}.  Proved as
\Cref{cor:ch20-semipositive-generic-nef}.

\cref{itm:ch20-impl:ntos}.  The first assertion is the composition of the
two previous ones.  For the second, let \(\rho:\mathcal Y\to\mathcal X\) be
a domination that is not an isomorphism and let \(Z\subseteq\mathcal Y_s\)
be a positive-dimensional fibre of \(\rho\); then
\((\rho^*\theta_{\mathcal X})^{\dim Z}\cdot Z=0\), so numerical positivity
fails on \(\mathcal Y\) however positive \(\theta_{\mathcal X}\) was.
Semipositivity is preserved and reflected by
\Cref{prop:ch20-semipositive-well-defined}.
\end{proof}

\begin{theorem}
\index{rigidity!of vertical nef classes}
\label{thm:ch20-vertical-nef-rigidity}
Assume \eqref{eq:ch20-standing} and let \(\mathcal X\) be a regular model.
Let \(D\in\Div_0(\mathcal X)_{\mathbf R}\).  The following are equivalent:
\begin{enumerate}
\item\label{itm:ch20-rig:nef} \([D]\) is relatively nef on \(\mathcal X\);
\item\label{itm:ch20-rig:zero} \([D]=0\) in \(N^1(\mathcal X/R)\);
\item\label{itm:ch20-rig:prop} \(D\in\mathbf R\,\mathcal X_s\);
\item\label{itm:ch20-rig:const} \(\varphi_D\) is constant.
\end{enumerate}
Consequently, for a model function \(\varphi\), the closed model form
\(\ddc\varphi\) is semipositive if and only if \(\varphi\) is constant; and
the only \(0\)-psh model functions are the constants.
\end{theorem}

\begin{proof}
\cref{itm:ch20-rig:zero}\(\Leftrightarrow\)\cref{itm:ch20-rig:prop}\(\Leftrightarrow\)\cref{itm:ch20-rig:const}
is \Cref{thm:ch20-ddc-exact} (exactness at
\(\Div_0(\mathcal X)_{\mathbf R}\)) together with
\eqref{eq:ch20-sign}, and
\cref{itm:ch20-rig:zero}\(\Rightarrow\)\cref{itm:ch20-rig:nef} is trivial.

Assume \cref{itm:ch20-rig:nef}.  Fix a relatively ample \(\mathcal A\) on
\(\mathcal X\) and write \(\mathcal X_s=\sum_ib_iE_i\).  Replacing
\(\mathcal A\) by a multiple, we may assume it relatively very ample; then
for each \(i\) the class \(\mathcal A^{n-1}\cdot[E_i]\) is represented by an
effective one-cycle \(\gamma_i\) on \(\mathcal X_s\), obtained by
intersecting \(E_i\) with \(n-1\) general members of the linear system.
Since \([D]\) is relatively nef,

\begin{equation}
\label{eq:ch20-rigidity-rows}
 (D\cdot E_i\cdot\mathcal A^{n-1})=(D\cdot\gamma_i)\geq0
 \qquad\text{for every }i .
\end{equation}

On the other hand \(\cO(\mathcal X_s)\simeq\cO_{\mathcal X}\) gives

\[
 0=(D\cdot\mathcal X_s\cdot\mathcal A^{n-1})
  =\sum_ib_i\,(D\cdot E_i\cdot\mathcal A^{n-1}),
\]

a sum of nonnegative terms with positive coefficients \(b_i\).  Hence every
term vanishes:
\((D\cdot E_i\cdot\mathcal A^{n-1})=0\) for all \(i\).  Writing
\(D=\sum_ia_iE_i\) and using linearity,

\[
 (D^2\cdot\mathcal A^{n-1})=\sum_ia_i\,(D\cdot E_i\cdot\mathcal A^{n-1})=0 .
\]

By \eqref{eq:ch20-hodge-laplacian} this forces
\(\sum_{i<j}w_{ij}(u_i-u_j)^2=0\) with \(u_i=a_i/b_i\), hence \(u_i=u_j\)
whenever \(E_i\cap E_j\neq\varnothing\); connectedness of the dual graph
gives \(u_i\equiv u\) and \(D=u\,\mathcal X_s\).  This is
\cref{itm:ch20-rig:prop}.

The final assertions follow: \(\ddc\varphi_D=[D]\) is semipositive if and
only if \([D]\) is relatively nef on some, equivalently any, determination
(\Cref{prop:ch20-semipositive-well-defined}), and by the above this happens
only for constant \(\varphi_D\).
\end{proof}

\begin{example}[Nef generic class, non-semipositive form]
\label{ex:ch20-nef-generic-not-semipositive}
Let \(\varphi\) be any nonconstant model function---for instance
\(\varphi=\varphi_E\) on \(\mathcal X_1=\operatorname{Bl}_p\mathbf P^2_R\)
in \Cref{ex:ch20-p2-blowup-numerics}, or
\(\varphi_D\) with \(D=AC_0+BC_1\), \(A\neq B\), on the running conic of
\Cref{ex:ch20-conic-ddc}.  Then \(\theta:=\ddc\varphi\) has
\(\{\theta\}=0\), which is nef, and \(\theta\) is not semipositive by
\Cref{thm:ch20-vertical-nef-rigidity}.  In the conic case one sees the
failure by hand: the degrees of \(\ddc\varphi\) on the two components are
\((B-A,A-B)\), which cannot both be nonnegative unless \(A=B\).

More is true, and it is worth stating because it recurs in
\Cref{chap:bfj-envelopes}: for an \emph{ample} generic class the same
phenomenon persists.  Take \(\theta_0=c_1(\cL)\) on the conic, with degrees
\((1,1)\), so that \(\{\theta_0\}\) is ample of degree \(2\), and put
\(\theta=\theta_0+\ddc\varphi_D\) with \(A=0,B=2\).  Then
\(\{\theta\}=\{\theta_0\}\) is ample, while the degrees of \(\theta\) are
\((3,-1)\), so \(\theta\) is not semipositive.
\Cref{prop:ch20-positive-representative} repairs this by adding
\(\psi=-\varphi_D\), at the cost of a shift of size \(2\) in sup norm; the
point of the smallness clause \eqref{eq:ch20-small-positive-representative}
is that no such shift is needed when \(\theta\) was semipositive to begin
with.
\end{example}

\begin{warning}{``Semipositive'' is a property of \(\theta\), not of a model}
Of the four notions in \Cref{def:ch20-four-positivities}, only
\cref{itm:ch20-pos:semi} and \cref{itm:ch20-pos:generic} are properties of the
closed model form.  Both \cref{itm:ch20-pos:ample} and
\cref{itm:ch20-pos:numerical} refer to a chosen model and are destroyed by
passing to a higher one.  In particular a phrase such as ``\(\theta\) is
ample'' has no meaning in \(\mathcal Z^{1,1}(X)\) and is never used in this
book; what is used is ``\(\theta\) is \(\mathcal X\)-positive'', with the
model displayed.  Statements in the literature about ``an ample model
form'' should be read as statements about a pair \((\mathcal X,\theta_{\mathcal X})\).
\end{warning}

\section{Intersection pairings and Hodge-type inequalities}
\label{sec:ch20-intersection}

\subsection{The vertical-factor \texorpdfstring{\((n+1)\)}{(n+1)}-fold pairing}
\label{sec:ch20-pairing}

A model \(\mathcal X\) has dimension \(n+1\), but relative numerical
equivalence controls degrees only on vertical curves.  Consequently the
intersection of \(n+1\) arbitrary classes in \(N^1(\mathcal X/R)\) does not
descend to closed model forms.  The intersection that does descend has an
actual vertical divisor in one slot and \(n\) relative numerical classes in
the others.  This is the pairing used in \cite[\S2.7]{BFJ15}.

\begin{proposition}
\index{intersection pairing!with a vertical divisor}
\label{prop:ch20-intersection-pairing}
There is a unique multilinear pairing, symmetric in its last \(n\) entries,

\begin{equation}
\label{eq:ch20-intersection-pairing}
 \mathcal D(X)_{\mathbf R}\times\mathcal Z^{1,1}(X)^n
 \longrightarrow\mathbf R,
 \qquad
 (\varphi_D,\theta_1,\ldots,\theta_n)\longmapsto
 (D\cdot\theta_1\cdots\theta_n),
\end{equation}

whose value when \(D\) and the \(\theta_i\) are determined on a common model
\(\mathcal X\) is

\[
 (D\cdot\theta_{1,\mathcal X}\cdots\theta_{n,\mathcal X})_{\mathcal X}
 =\sum_E\ord_E(D)
   (\theta_{1,\mathcal X}|_E\cdots\theta_{n,\mathcal X}|_E),
\]

where \(E\) ranges over the irreducible components of \(\mathcal X_s\).
\end{proposition}

\begin{proof}
Existence and uniqueness reduce to independence of the common determination.
Let \(\rho:\mathcal Y\to\mathcal X\) be a domination.  Since \(\rho\) is
birational and proper, \(\rho_*[\mathcal Y]=[\mathcal X]\), so the
projection formula for Chern classes gives

\[
 (\rho^*D\cdot\rho^*\alpha_1\cdots\rho^*\alpha_n)_{\mathcal Y}
 =(D\cdot\alpha_1\cdots\alpha_n)_{\mathcal X}.
\]

This is the Chow projection formula \Cref{thm:appF-projection} applied to
the vertical cycle \(D\), and is also the case \(d=1\) of
\Cref{thm:model-projection-formula}.  If two collections of data represent
the same model function and closed model forms, pass to a common domination
by \Cref{prop:common-domination}.  The divisor pulls back because it
represents the same model function, the relative classes agree after
pullback by \Cref{prop:ch20-two-model-comparison}, and the displayed identity
shows that the two values agree.  Multilinearity and symmetry in the
\(\theta_i\) are \Cref{prop:appF-mixed-cartier}.  Uniqueness follows because
all finitely many entries admit a common determination
(\eqref{eq:ch20-model-functions} and
\Cref{cor:ch20-limit-structure}).
\end{proof}

The first slot in \eqref{eq:ch20-intersection-pairing} remembers the divisor
\(D\), equivalently the model function \(\varphi_D\), and not merely its
relative class \([D]=\ddc\varphi_D\).  This distinction is essential: for
\(D=\mathcal X_s\), one has \(\varphi_D=1\) but
\(\ddc\varphi_D=0\).  The model-function pairing of
\Cref{prop:metrics-cl-model-function-integral} is

\begin{equation}
\label{eq:ch20-pairing-as-integral}
 (D\cdot\theta_1\cdots\theta_n)
 =\int_{X^{\An}}\varphi_D\;\MA(\theta_1,\ldots,\theta_n),
\end{equation}

where \(\MA(\theta_1,\ldots,\theta_n)\) is the mixed Chambert--Loir measure
\eqref{eq:appJ-cl-measure} of the corresponding model metrics---a signed
atomic measure in general, positive when the \(\theta_i\) are semipositive
(\Cref{prop:ch20-measure-positivity}).  Taking \(D=\mathcal X_s\), so
\(\varphi_D=1\), gives the total mass
\((\{\theta_1\}\cdots\{\theta_n\})\), which is
\eqref{eq:ch20-special-fibre-degree}.  Equation
\eqref{eq:ch20-pairing-as-integral} is the bridge between the algebra of this
chapter and the measure theory of \Cref{chap:bfj-ma}; the following
corollary is the model case of the integration-by-parts formula
\Cref{prop:model-integration-by-parts} proved there.

\begin{corollary}
\label{cor:ch20-integration-by-parts}
For model functions \(f=\varphi_D\), \(g=\varphi_{D'}\) and closed model
forms \(\eta_1,\ldots,\eta_{n-1}\),

\begin{equation}
\label{eq:ch20-ibp}
\begin{aligned}
 \int_{X^{\An}}f\,\MA(\ddc g,\eta_1,\ldots,\eta_{n-1})
 &=(D\cdot D'\cdot\eta_1\cdots\eta_{n-1})\\
 &=\int_{X^{\An}}g\,\MA(\ddc f,\eta_1,\ldots,\eta_{n-1}).
\end{aligned}
\end{equation}
\end{corollary}

\begin{proof}
Choose a common determination and apply
\eqref{eq:ch20-pairing-as-integral} twice; the middle term is symmetric in
\(D\) and \(D'\) by \Cref{prop:model-intersection-basic}.
\end{proof}

\subsection{The Hodge form}
\label{sec:ch20-hodge}

Choose a regular model \(\mathcal X\), write
\eqref{eq:ch20-special-fibre-components}, and fix a relatively ample
\(\mathcal A\).  If \(D=\sum_i a_iE_i\) and
\(F=\sum_i c_iE_i\), define

\begin{equation}
\label{eq:ch20-dirichlet-pairing}
             \langle D,F\rangle_{\mathcal A}
             :=-\bigl(D\cdot F\cdot\mathcal A^{n-1}\bigr).
\end{equation}

The calculation in the proof of \Cref{thm:ch20-ddc-exact} yields the
following precise form of the relative Hodge inequality.

\begin{proposition}\label{prop:ch20-relative-hodge}
For every vertical real Cartier divisor \(D\),

\begin{equation}
\label{eq:ch20-relative-hodge}
                    (D^2\cdot\mathcal A^{n-1})\leq0.
\end{equation}

Equality holds exactly when \(D\) is numerically proportional to
\(\mathcal X_s\).  Hence \eqref{eq:ch20-dirichlet-pairing} descends to an
inner product on vertical divisors modulo constants, and

\begin{equation}
\label{eq:ch20-hodge-cauchy-schwarz}
 |D\cdot F\cdot\mathcal A^{n-1}|^2
 \leq
 (D^2\cdot\mathcal A^{n-1})
 (F^2\cdot\mathcal A^{n-1}).
\end{equation}
\end{proposition}

\begin{proof}
Set \(q(D)=\langle D,D\rangle_{\mathcal A}\).  The relative Hodge-index
calculation in the proof of \Cref{thm:ch20-ddc-exact} shows that \(q\) is
positive semidefinite and that its nullspace is
\(\mathbf R\mathcal X_s\).  This is precisely
\eqref{eq:ch20-relative-hodge}, including its equality statement.  Since
\(\langle\cdot,\cdot\rangle_{\mathcal A}\) is symmetric and vanishes on
\(\mathbf R\mathcal X_s\), it therefore induces a positive-definite
bilinear form on the quotient by \(\mathbf R\mathcal X_s\).  Applying the
ordinary Cauchy--Schwarz inequality in that quotient gives

\[
 \bigl|-(D\cdot F\cdot\mathcal A^{n-1})\bigr|^2
 \leq
 \bigl(-(D^2\cdot\mathcal A^{n-1})\bigr)
 \bigl(-(F^2\cdot\mathcal A^{n-1})\bigr),
\]

which is exactly \eqref{eq:ch20-hodge-cauchy-schwarz} after cancelling the
two minus signs on the right.
\end{proof}

Both factors on the right of \eqref{eq:ch20-hodge-cauchy-schwarz} are
nonpositive, so their product is nonnegative.  The sign in
\eqref{eq:ch20-relative-hodge} is the algebraic origin of concavity of the
Monge--Amp\`ere energy in \Cref{chap:bfj-energy}.

The version actually needed downstream replaces the powers of a single
relatively ample class by an arbitrary tuple of relatively nef classes.  It
is proved by exactly the same computation, followed by a limit.

\begin{theorem}
\index{Hodge index theorem!relative, mixed form}
\label{thm:ch20-mixed-hodge}
Assume \eqref{eq:ch20-standing}, let \(\mathcal X\) be a regular model, let
\(\alpha_1,\ldots,\alpha_{n-1}\in N^1(\mathcal X/R)\) be relatively nef, and
let \(D\in\Div_0(\mathcal X)_{\mathbf R}\).  Then

\begin{equation}
\label{eq:ch20-mixed-hodge}
 \bigl(D^2\cdot\alpha_1\cdots\alpha_{n-1}\bigr)\leq0 ,
\end{equation}

and for \(D,F\in\Div_0(\mathcal X)_{\mathbf R}\),

\begin{equation}
\label{eq:ch20-mixed-cauchy-schwarz}
 \bigl|(D\cdot F\cdot\alpha_1\cdots\alpha_{n-1})\bigr|^2
 \leq
 \bigl(D^2\cdot\alpha_1\cdots\alpha_{n-1}\bigr)
 \bigl(F^2\cdot\alpha_1\cdots\alpha_{n-1}\bigr).
\end{equation}

If moreover all \(\alpha_i\) are relatively ample, equality holds in
\eqref{eq:ch20-mixed-hodge} if and only if \(D\in\mathbf R\,\mathcal X_s\).
\end{theorem}

\begin{proof}
\emph{Step 1: the ample case.}  Assume first that every \(\alpha_i\) is
relatively ample.  Write \(\mathcal X_s=\sum_ib_iE_i\), \(D=\sum_ia_iE_i\),
\(u_i=a_i/b_i\), and set

\[
 c_{ij}:=(E_i\cdot E_j\cdot\alpha_1\cdots\alpha_{n-1}),
 \qquad
 w_{ij}:=b_ib_jc_{ij}\quad(i\ne j).
\]

For \(i\ne j\), the cycle \(E_i\cdot E_j\) is effective of dimension
\(n-1\), being a proper intersection of two distinct prime divisors on the
regular scheme \(\mathcal X\); a product of \(n-1\) ample classes has
nonnegative degree on an effective \((n-1)\)-cycle, so \(w_{ij}\geq0\).
Moreover \(\cO(\mathcal X_s)\simeq\cO_{\mathcal X}\) gives the row-sum
relation \(\sum_jb_jc_{ij}=0\) for every \(i\).  The computation displayed
in Step~1 of the proof of \Cref{thm:ch20-ddc-exact} used only these two
facts, so it applies verbatim and yields

\begin{equation}
\label{eq:ch20-mixed-hodge-laplacian}
 -\bigl(D^2\cdot\alpha_1\cdots\alpha_{n-1}\bigr)
 =\sum_{i<j}w_{ij}(u_i-u_j)^2\;\geq\;0,
\end{equation}

which is \eqref{eq:ch20-mixed-hodge}.  Equality forces \(u_i=u_j\) whenever
\(w_{ij}>0\).  For ample \(\alpha_i\) one has \(w_{ij}>0\) as soon as
\(E_i\cap E_j\neq\varnothing\), because then \(E_i\cdot E_j\) is a nonzero
effective cycle and ample classes have strictly positive degree on it;
connectedness of the dual graph of \(\mathcal X_s\), established in Step~1 of
\Cref{thm:ch20-ddc-exact}, then forces all \(u_i\) equal, that is
\(D\in\mathbf R\mathcal X_s\).

\emph{Step 2: the nef case by approximation.}  Fix a relatively ample
\(\mathcal A\) and put \(\alpha_i^\varepsilon=\alpha_i+\varepsilon\mathcal A\),
which is relatively ample for \(\varepsilon>0\) by
\Cref{lem:ch20-special-fibre-test}\,\cref{itm:ch20-sft:cone}.  Step~1 gives
\((D^2\cdot\alpha_1^\varepsilon\cdots\alpha_{n-1}^\varepsilon)\leq0\), and
the left-hand side is a polynomial in \(\varepsilon\); letting
\(\varepsilon\to0^+\) gives \eqref{eq:ch20-mixed-hodge}.

\emph{Step 3: Cauchy--Schwarz.}  The symmetric bilinear form
\(\langle D,F\rangle:=-(D\cdot F\cdot\alpha_1\cdots\alpha_{n-1})\) on
\(\Div_0(\mathcal X)_{\mathbf R}\) is positive semidefinite by
\eqref{eq:ch20-mixed-hodge}.  The Cauchy--Schwarz inequality is valid for
any positive semidefinite symmetric bilinear form on a real vector space---no
definiteness is needed---so
\(|\langle D,F\rangle|^2\leq\langle D,D\rangle\langle F,F\rangle\), which is
\eqref{eq:ch20-mixed-cauchy-schwarz} after cancelling the four minus signs.
\end{proof}

\begin{remark}[Where the two forms of the inequality are used]
\label{rem:ch20-hodge-uses}
\Cref{prop:ch20-relative-hodge} is the case
\(\alpha_1=\cdots=\alpha_{n-1}=\mathcal A\); it is what
\Cref{thm:ch20-ddc-exact} needs, because there the equality case has to be
sharp.  \Cref{thm:ch20-mixed-hodge} is what
\Cref{chap:bfj-ma,chap:bfj-energy,chap:bfj-capacity} need, because there the
classes \(\alpha_i\) are curvature forms \(\theta+\ddc\varphi_i\) of
different psh potentials and are only nef.  In the notation of
\Cref{prop:model-integration-by-parts}, \eqref{eq:ch20-mixed-hodge} is
exactly \eqref{eq:model-dirichlet-negativity}, and
\eqref{eq:ch20-mixed-cauchy-schwarz} is the Cauchy--Schwarz inequality
invoked in \Cref{cor:bounded-ma-multilinearity}, from which the capacity
estimates of \Cref{chap:bfj-capacity} are derived.  We emphasise that no
equality case is claimed in the merely nef situation: for
\(\alpha_1=\cdots=\alpha_{n-1}=0\) the left-hand side of
\eqref{eq:ch20-mixed-hodge} vanishes for every \(D\).
\end{remark}

\begin{corollary}
\label{cor:ch20-concavity}
Let \(\theta_0\in\mathcal Z^{1,1}(X)\) be semipositive, let \(\varphi\) be a
model function, and suppose that \(\theta_t:=\theta_0+t\,\ddc\varphi\) is
semipositive for every \(t\) in an interval \(J\subseteq\mathbf R\).  Then
the mixed vertical energy difference

\begin{equation}
\label{eq:ch20-concavity}
 t\longmapsto E_{\theta_0}(t\varphi)
 :=\frac1{n+1}\sum_{j=0}^n
   \bigl(tD\cdot\theta_t^{\,j}\cdot\theta_0^{\,n-j}\bigr),
\end{equation}

where \(D\) represents \(\varphi\) on a common model, is concave on \(J\).
More precisely,

\[
 \frac{\rd}{\rd t}E_{\theta_0}(t\varphi)
   =(D\cdot\theta_t^{\,n})
   =\int_{X^{\An}}\varphi\,\MA(\theta_t^{\,n}),
 \qquad
 \frac{\rd^2}{\rd t^2}E_{\theta_0}(t\varphi)
   =n(D^2\cdot\theta_t^{\,n-1})\leq0.
\]
\end{corollary}

\begin{proof}
Choose a model \(\mathcal X\) determining \(\theta_0\) and
\(\varphi=\varphi_D\).  Every summand in
\eqref{eq:ch20-concavity} is defined by
\eqref{eq:ch20-intersection-pairing}, since its first factor is the vertical
divisor \(tD\).  Put
\(A_j(t):=(D\cdot\theta_t^j\cdot\theta_0^{n-j})\).  Multilinearity and
\(t[D]=\theta_t-\theta_0\) give, for \(j\geq1\),

\[
 t(D^2\cdot\theta_t^{j-1}\cdot\theta_0^{n-j})
 =A_j(t)-A_{j-1}(t).
\]

Differentiating the finite sum in \eqref{eq:ch20-concavity} and using this
identity makes all terms telescope, leaving
\(\frac{\rd}{\rd t}E_{\theta_0}(t\varphi)=A_n(t)\).  A second derivative
gives \(n(D^2\cdot\theta_t^{n-1})\).  The integral identity is
\eqref{eq:ch20-pairing-as-integral}.  For \(t\in J\) the class
\(\theta_{t,\mathcal X}\) is relatively nef by
\Cref{prop:ch20-semipositive-well-defined}, so
\Cref{thm:ch20-mixed-hodge} supplies the final inequality.
\end{proof}

\Cref{cor:ch20-concavity} is the model-level shadow of the concavity of the
Monge--Amp\`ere energy proved in \Cref{chap:bfj-energy}: the mixed vertical
sum in \eqref{eq:ch20-concavity} is exactly the model case of the
normalization \eqref{eq:appJ-energy}.  Thus the concavity statement uses only
well-defined vertical-factor intersections; no absolute
\((n+1)\)-fold self-intersection of a closed model form is invoked.

\subsection{Componentwise positivity and mixed degrees}
\label{sec:ch20-componentwise}

If \(\alpha_1,\ldots,\alpha_n\) are relatively nef classes, then their
restrictions to each component \(E_i\) are nef, and therefore

\begin{equation}
\label{eq:ch20-mixed-component-positive}
        (\alpha_1|_{E_i}\cdots\alpha_n|_{E_i})\geq0.
\end{equation}

Summing with the multiplicities \(b_i\) and using
\eqref{eq:ch20-special-fibre-degree} recovers the generic mixed degree.
The Khovanskii--Teissier inequalities on the projective varieties \(E_i\)
give log-concavity of the corresponding mixed sequences.  We use only these
standard componentwise inequalities; no undefined intersection of two forms
living on different models is taken.

\begin{proposition}
\label{prop:ch20-measure-positivity}
Let \(\theta_1,\ldots,\theta_n\in\mathcal Z^{1,1}(X)\) be semipositive and
determined on a common model \(\mathcal X\).  Then the mixed model measure

\begin{equation}
\label{eq:ch20-model-measure}
 \MA(\theta_1,\ldots,\theta_n)
 =\sum_{i\in I}b_i\,
   \bigl(\theta_{1}|_{E_i}\cdots\theta_{n}|_{E_i}\bigr)\,
   \delta_{x_{E_i}}
\end{equation}

is a positive atomic measure supported on the divisorial points of
\(\mathcal X\), of total mass \((\{\theta_1\}\cdots\{\theta_n\})\).  For a
model function \(\varphi_D\),

\begin{equation}
\label{eq:ch20-measure-pairing}
 \int_{X^{\An}}\varphi_D\,\MA(\theta_1,\ldots,\theta_n)
 =(D\cdot\theta_1\cdots\theta_n)_{\mathcal X}.
\end{equation}
\end{proposition}

\begin{proof}
Positivity of the atoms is \eqref{eq:ch20-mixed-component-positive}, using
that the restriction of a relatively nef class to a component of
\(\mathcal X_s\) is nef and that a product of nef classes on a projective
scheme over a field has nonnegative degree.  Independence of the model,
symmetry and the total mass are \Cref{thm:cl-model-properties}, and
\eqref{eq:ch20-measure-pairing} is
\Cref{prop:metrics-cl-model-function-integral}, that is,
\eqref{eq:metrics-cl-pairing}.  The total mass is
\eqref{eq:ch20-special-fibre-degree}, obtained from
\eqref{eq:ch20-measure-pairing} by taking \(D=\mathcal X_s\) and
\(\varphi_D=1\).
\end{proof}

\begin{warning}{Log-concavity is componentwise, not global}
The Khovanskii--Teissier inequalities quoted after
\eqref{eq:ch20-mixed-component-positive} hold on each projective
\(k\)-variety \(E_i\) separately.  They do \emph{not} directly give
log-concavity of the sequence of mixed masses
\(\int\MA(\theta^{(j)},\theta'^{(n-j)})\), because that sequence is a
\(b_i\)-weighted sum of componentwise sequences, and a sum of log-concave
sequences need not be log-concave.  What is true, and what
\Cref{chap:bfj-energy} uses, is the Hodge-type inequality
\eqref{eq:ch20-mixed-cauchy-schwarz}, which is a statement about
\emph{vertical} classes and follows from the row-sum relation rather than
from componentwise positivity.  The two mechanisms are independent and
should not be conflated.
\end{warning}

\section{Worked examples}
\label{sec:ch20-examples}

The five examples below are carried to explicit numbers.  Together they test
every normalization in the chapter: the sign of \(\varphi_D\), the sign of
\(\ddc\), the degrees defining relative nefness, the direct-limit
identification, and both Hodge inequalities.

\subsection{A semistable curve and its Laplacian}
\label{sec:ch20-ex-curve}

\begin{example}[A semistable curve]
\label{ex:ch20-semistable-curve}
Let \(n=1\) and let \(\mathcal X\) be a regular semistable model, initially
allowing multiplicities \(b_i\).  Put

\begin{equation}
\label{eq:ch20-curve-values}
       D=\sum_i a_iE_i,
       \qquad u_i=\varphi_D(x_{E_i})=\frac{a_i}{b_i}.
\end{equation}

Since \((\mathcal X_s\cdot E_i)=0\),

\begin{equation}
\label{eq:ch20-graph-laplacian}
 \begin{aligned}
 (D\cdot E_i)
 &=\sum_j b_ju_j(E_j\cdot E_i)\\
 &=\sum_{j\neq i}b_j(E_j\cdot E_i)(u_j-u_i).
 \end{aligned}
\end{equation}

Thus the degree of \(D\) on \(E_i\) is a weighted graph Laplacian with the
sign dictated by \eqref{eq:ch20-sign}.  If \(\theta\) has component degrees
\(d_i=(\theta\cdot E_i)\), then \(\theta+\ddc\varphi_D\) is nef exactly
when

\begin{equation}
\label{eq:ch20-curve-nef-test}
 d_i+\sum_{j\neq i}b_j(E_j\cdot E_i)(u_j-u_i)\geq0
 \quad\text{for every }i.
\end{equation}

For a reduced two-component fibre with one node, the intersection matrix is

\begin{equation}
\label{eq:ch20-two-component-matrix}
              \begin{pmatrix}-1&1\\1&-1\end{pmatrix}.
\end{equation}

Consequently
\((D^2)=-(u_1-u_2)^2\), visibly confirming
\eqref{eq:ch20-relative-hodge}; equality means that the two vertex values
are equal, hence that the model function is constant.
\end{example}

\begin{example}[A chain of four components, with numbers]
\label{ex:ch20-chain-three}
Let \(X\) be a smooth projective curve admitting a regular semistable model
\(\mathcal X\) whose special fibre is a chain
\(E_0,E_1,E_2,E_3\) of four components, all of multiplicity one, with
\((E_{i}\cdot E_{i+1})=1\) and \((E_i\cdot E_j)=0\) for \(|i-j|\geq2\); the
resolution of the \(A_2\) singularity \(xy=\pi^3\) in
\Cref{ex:ch08-An-edge} produces such a chain locally.  The row-sum relation
\((\mathcal X_s\cdot E_i)=0\) then gives \((E_0^2)=(E_3^2)=-1\) and
\((E_1^2)=(E_2^2)=-2\).  The dual complex is the segment \([0,3]\)
subdivided at the integers, and a model function determined on this model is
the piecewise affine function with vertex values \(u_0,u_1,u_2,u_3\).

Let \(\theta\) be a closed model form with degrees
\((d_0,d_1,d_2,d_3)=(1,0,0,1)\); its generic degree is
\(\sum_id_i=2\).  By \eqref{eq:ch20-curve-nef-test}, \(\varphi_D\) is
\(\theta\)-psh precisely when

\begin{equation}
\label{eq:ch20-chain-constraints}
 1+(u_1-u_0)\geq0,\quad
 u_0+u_2-2u_1\geq0,\quad
 u_1+u_3-2u_2\geq0,\quad
 1+(u_2-u_3)\geq0 .
\end{equation}

The two middle inequalities say that \(u\) is \emph{convex} at the two
interior vertices---which is the correct sign, since by
\Cref{prop:ch20-curve-laplacian} a convex corner has positive
Laplacian---and the two outer ones are one-sided slope bounds at the ends.
Normalizing \(\max_iu_i=0\), two extreme admissible vectors are:

\begin{itemize}
\item \(u=(0,0,0,0)\), for which \(\theta+\ddc\varphi=\theta\) and the
 measure is \(\delta_{x_{E_0}}+\delta_{x_{E_3}}\);
\item \(u=(0,-1,-1,0)\), a convex tent, for which the four left-hand sides of
 \eqref{eq:ch20-chain-constraints} are \(0,1,1,0\), so the degrees of
 \(\theta+\ddc\varphi\) are \((0,1,1,0)\) and the measure is
 \(\delta_{x_{E_1}}+\delta_{x_{E_2}}\).
\end{itemize}

In both cases the total mass is \(2\), as
\eqref{eq:ch20-special-fibre-degree} demands: the perturbation
\(\ddc\varphi\) has total degree zero.  Note how the mass has been
transported from the two ends of the segment to its two interior vertices;
this is the model-level version of \Cref{ex:edge-mass-transport}.

Finally, by \eqref{eq:ch20-hodge-laplacian} with \(n=1\) and all edge
weights equal to \(1\), the Hodge form is
\((D^2)=-\sum_{i=0}^{2}(u_i-u_{i+1})^2\); for the tent this is
\(-\bigl(1+0+1\bigr)=-2<0\), confirming
\eqref{eq:ch20-relative-hodge}, and it vanishes exactly for constant \(u\),
confirming the equality case.
\end{example}

\subsection{Two models of a surface and the direct limit}
\label{sec:ch20-ex-two-model}

\begin{example}[Two models of \(\mathbf P^2\), and what the limit identifies]
\label{ex:ch20-two-model-surface}
We continue \Cref{ex:ch20-p2-blowup-numerics}, now with two blowups.  Let
\(X=\mathbf P^2_K\), \(\cX_0=\mathbf P^2_R\), \(F=\cX_{0,s}\), and let
\(p\neq q\) be two distinct \(k\)-rational points of \(F\).  Put

\[
 \cX_1=\operatorname{Bl}_p\cX_0,\qquad
 \cX_2=\operatorname{Bl}_q\cX_0,\qquad
 \cX_3=\operatorname{Bl}_{\{p,q\}}\cX_0 ,
\]

with dominations \(\rho_i:\cX_i\to\cX_0\) for \(i=1,2\), and
\(\sigma_i:\cX_3\to\cX_i\), \(\pi=\rho_i\circ\sigma_i:\cX_3\to\cX_0\).
Neither \(\cX_1\) nor \(\cX_2\) dominates the other: the birational map
between them contracts \(E_p\) and extracts \(E_q\), so it is not a
morphism in either direction.  They are the smallest pair of incomparable
models of a surface.

\emph{The three numerical spaces.}  Write \(E_p,E_q\) for the exceptional
divisors on \(\cX_1,\cX_2\) and \(\widehat E_p=\sigma_1^*E_p\),
\(\widehat E_q=\sigma_2^*E_q\) on \(\cX_3\); write \(F_i\) for the strict
transform of \(F\) on \(\cX_i\).  Then

\[
 \cX_{1,s}=F_1+E_p,\qquad
 \cX_{2,s}=F_2+E_q,\qquad
 \cX_{3,s}=F_3+\widehat E_p+\widehat E_q ,
\]

all multiplicities being one, and by \Cref{cor:ch20-dimension-count},

\begin{equation}
\label{eq:ch20-two-model-dimensions}
 \dim N^1(\cX_0/R)=1,\quad
 \dim N^1(\cX_1/R)=\dim N^1(\cX_2/R)=2,\quad
 \dim N^1(\cX_3/R)=3 ,
\end{equation}

with bases \(\{h\}\), \(\{\rho_1^*h,[E_p]\}\), \(\{\rho_2^*h,[E_q]\}\) and
\(\{\pi^*h,[\widehat E_p],[\widehat E_q]\}\) respectively, by the argument
of \Cref{ex:ch20-p2-blowup-numerics} applied at each point.

\emph{The two images inside \(N^1(\cX_3/R)\).}  The maps \(\sigma_i^*\) are
injective (\Cref{lem:ch20-pullback-injective}) and send

\[
 \sigma_1^*(\rho_1^*h)=\pi^*h,\quad
 \sigma_1^*[E_p]=[\widehat E_p];
 \qquad
 \sigma_2^*(\rho_2^*h)=\pi^*h,\quad
 \sigma_2^*[E_q]=[\widehat E_q].
\]

Hence, inside the three-dimensional space \(N^1(\cX_3/R)\),

\begin{equation}
\label{eq:ch20-two-model-images}
 \begin{aligned}
 \sigma_1^*N^1(\cX_1/R)
   &=\langle\pi^*h,[\widehat E_p]\rangle,\\
 \sigma_2^*N^1(\cX_2/R)
   &=\langle\pi^*h,[\widehat E_q]\rangle,\\
 \sigma_1^*N^1(\cX_1/R)\cap\sigma_2^*N^1(\cX_2/R)
   &=\langle\pi^*h\rangle .
 \end{aligned}
\end{equation}

The intersection of the two images is exactly the image of
\(N^1(\cX_0/R)\).  So the direct limit identifies a class on \(\cX_1\) with
a class on \(\cX_2\) if and only if both are pulled back from \(\cX_0\),
and \emph{no other identification occurs}.

\emph{An identification that is not obvious from the coefficients.}
Consider

\[
\begin{aligned}
 \alpha_1&:=c_1(\rho_1^*\cO(1))+[F_1]+[E_p]&&\in N^1(\cX_1/R),\\
 \alpha_2&:=c_1(\rho_2^*\cO(1))+2[F_2]+2[E_q]&&\in N^1(\cX_2/R).
\end{aligned}
\]

Written in coordinates these look entirely different: the first has vertical
part \((1,1)\) in the basis \((F_1,E_p)\), the second has vertical part
\((2,2)\) in the basis \((F_2,E_q)\), and the two models do not even have
the same set of components.  But \([\cX_{i,s}]=0\), so
\(\alpha_1=\rho_1^*h\) and \(\alpha_2=\rho_2^*h\), and both pull back to
\(\pi^*h\) on \(\cX_3\).  By
\Cref{prop:ch20-two-model-comparison} they define the same closed model
form, and no further model needs to be inspected.  Note that the two
\emph{model metrics} attached to the two representatives differ: their
weights differ by the constant \(1\) and \(2\) respectively, since
\(\varphi_{\cX_{i,s}}=1\).  This is precisely the ambiguity recorded by the
first term of the exact sequence \eqref{eq:ch20-ddc-exact}.

\emph{An apparent identification that fails.}  Consider instead
\(\beta_1:=c_1(\rho_1^*\cO(1))+[E_p]\) and
\(\beta_2:=c_1(\rho_2^*\cO(1))+[E_q]\).  Both have generic class \(h\), and
both are ``the reference class corrected by one exceptional divisor''.  They
are \emph{not} equal in \(\mathcal Z^{1,1}(X)\): their pullbacks to
\(\cX_3\) are \(\pi^*h+[\widehat E_p]\) and \(\pi^*h+[\widehat E_q]\),
which differ because \([\widehat E_p]\neq[\widehat E_q]\) in
\(N^1(\cX_3/R)\); indeed \(([\widehat E_p]\cdot\ell)=-1\) and
\(([\widehat E_q]\cdot\ell)=0\) for a line \(\ell\subseteq\widehat E_p\).
Equivalently, the corresponding model functions
\(\varphi_{E_p}\) and \(\varphi_{E_q}\) are different functions on
\(X^{\An}\): the first takes the value \(1\) at \(x_{E_p}\) and \(0\) at
\(x_{E_q}\), the second the other way round.  The moral is the one stated
after \Cref{def:ch20-closed-model-form}: equality of closed model forms is
never equality of coefficient vectors, and it is always a computation on one
common domination.
\end{example}

\subsection{A vertical pairing with the Hodge index verified}
\label{sec:ch20-ex-hodge}

\begin{example}[Vertical triple products on a blown-up plane]
\label{ex:ch20-triple-products}
Keep \(\cX_1=\operatorname{Bl}_p\mathbf P^2_R\) of
\Cref{ex:ch20-p2-blowup-numerics}, so \(n=2\) and \(\dim\cX_1=3\); the
pairing \eqref{eq:ch20-intersection-pairing} has one vertical-divisor factor
and two relative-class factors.  We compute it using the space
\(N^1(\cX_1/R)\) spanned by \(\rho^*h\) and \([E]\).

\emph{The three basic vertical triple products.}
\begin{align}
\label{eq:ch20-triple-basic}
 \bigl(E\cdot(\rho^*h)^2\bigr)&=0, &
 \bigl(E^2\cdot\rho^*h\bigr)&=0, &
 \bigl(E^3\bigr)&=1 .
\end{align}
The first two vanish because \(\rho^*h|_E\) is trivial, \(E\) being
contracted to a point.  For the last,
\(\cO(E)|_E\simeq\cO_{\mathbf P^2}(-1)\) gives
\((E^3)=\bigl((\cO_{\mathbf P^2}(-1))^2\bigr)_{\mathbf P^2}=1\).
There is deliberately no entry \((\rho^*h)^3\): an absolute triple product
of relative numerical classes is not part of
\eqref{eq:ch20-intersection-pairing}.

\emph{A family of nef classes and its energy difference.}  Put
\(\theta_t:=2\rho^*h-t[E]\) for \(t\in\mathbf R\).  Restricting,
\(\theta_t|_E=\cO_{\mathbf P^2}(t)\) and
\(\theta_t|_{F_1}=2h_1-te\), so by
\Cref{lem:ch20-special-fibre-test} the class \(\theta_t\) is relatively nef
exactly for \(0\leq t\leq2\) and relatively ample exactly for
\(0<t<2\).  Expanding with \eqref{eq:ch20-triple-basic},

\begin{equation}
\label{eq:ch20-triple-energy}
 3E_{\theta_0}(t\varphi)
 =\sum_{j=0}^2\bigl(-tE\cdot\theta_t^{\,j}\cdot
                      \theta_0^{\,2-j}\bigr)
 =-t^3 .
\end{equation}

Here \(\theta_0=2\rho^*h\) and \(\varphi=-\varphi_E\), so each summand has
the required vertical first factor.  On \([0,2]\) the right-hand side is
concave, as \Cref{cor:ch20-concavity} predicts, with second derivative
\(-6t\leq0\); note that \(-t^3\) is \emph{convex} for \(t<0\), which is
consistent, because \(\theta_t\) is not semipositive there.  At \(t=1\),
\eqref{eq:ch20-triple-energy} gives
\(3E_{\theta_0}(\varphi)=-1\), consistent with \(\varphi\leq0\) and
\(E_{\theta_0}(0)=0\).

\emph{The Hodge inequality, checked.}  Take \(\mathcal A=\theta_1=2\rho^*h-[E]\),
which is relatively ample, and \(D=aE\) for \(a\in\mathbf R\).  Then

\begin{equation}
\label{eq:ch20-triple-hodge-check}
 \bigl(D^2\cdot\mathcal A\bigr)
 =a^2\bigl(E^2\cdot(2\rho^*h-E)\bigr)
 =a^2\bigl(2\rho^*hE^2-E^3\bigr)
 =-a^2\;\leq\;0 ,
\end{equation}

with equality if and only if \(a=0\), that is \(D=0\in\mathbf R\cX_{1,s}\).
This is \eqref{eq:ch20-relative-hodge} with its equality case.  The same
number is computed from \eqref{eq:ch20-hodge-laplacian}: the components are
\(F_1,E\) with \(b_{F_1}=b_E=1\), the vertex values of \(D=aE\) are
\(u_{F_1}=0\), \(u_E=a\), and the single edge weight is

\[
 w_{F_1E}=(F_1\cdot E\cdot\mathcal A)=(\mathcal A\cdot e)
 =\bigl(\cO_{\mathbf P^2}(1)\cdot(\text{line})\bigr)=1 ,
\]

so \(q(D)=w_{F_1E}(u_{F_1}-u_E)^2=a^2\), in agreement with
\eqref{eq:ch20-triple-hodge-check}.  Two independent computations of the
same number, one by expanding the vertical triple product and one by the
Laplacian formula,
agreeing---this is the smallest simultaneous test of
\eqref{eq:ch20-hodge-laplacian}, \eqref{eq:ch20-relative-hodge}, and the
projection formula.

\emph{Cauchy--Schwarz on three components.}  Pass to
\(\cX_3=\operatorname{Bl}_{\{p,q\}}\mathbf P^2_R\) of
\Cref{ex:ch20-two-model-surface} and take
\(\mathcal A_3=3\pi^*h-[\widehat E_p]-[\widehat E_q]\), relatively ample by
the same Nakai--Moishezon check as before.  With
\(D=a\widehat E_p\) and \(F=b\widehat E_q\),

\[
 (D^2\cdot\mathcal A_3)=-a^2,\qquad
 (F^2\cdot\mathcal A_3)=-b^2,\qquad
 (D\cdot F\cdot\mathcal A_3)=0 ,
\]

the last because \(\widehat E_p\) and \(\widehat E_q\) are disjoint.  Thus
\eqref{eq:ch20-hodge-cauchy-schwarz} reads \(0\leq a^2b^2\), a strict
inequality unless \(ab=0\): the two exceptional directions are orthogonal
for the Dirichlet pairing.  Equality is attained on proportional divisors,
for instance \(F=2D\), where both sides equal \(4a^4\).  Finally, the
Laplacian formula gives
\(q(a\widehat E_p+b\widehat E_q)=w_{F_3\widehat E_p}a^2
 +w_{F_3\widehat E_q}b^2+w_{\widehat E_p\widehat E_q}(a-b)^2
 =a^2+b^2+0\), matching
\(-\bigl((a\widehat E_p+b\widehat E_q)^2\cdot\mathcal A_3\bigr)=a^2+b^2\).
\end{example}

\subsection{Maxima and the toric picture}
\label{sec:ch20-ex-max-toric}

\begin{example}[A maximum on the nodal interval]
\label{ex:ch20-nodal-max}
On the local model \(xy=\pi\), parameterize the skeleton by
\(t=v(x)\in[0,1]\), so \(v(y)=1-t\).  Let

\[
              \mathfrak a=(x^p,y^q),\qquad p,q\geq1.
\]

Then

\begin{equation}
\label{eq:ch20-nodal-ideal}
 \log|\mathfrak a|(t)
   =\max\{-pt,-q(1-t)\}
   =-\min\{pt,q(1-t)\}.
\end{equation}

The two affine branches meet at \(t=q/(p+q)\).  Formula
\eqref{eq:ch20-nodal-ideal} is exactly the ideal-theoretic maximum in
\eqref{eq:ch20-max-ideal-function}; the normalized blowup inserts the
breakpoint when it is rational.

Two numerical remarks.  First, the identity
\eqref{eq:ch20-nodal-ideal} is a special case of
\eqref{eq:ch20-log-ideal}: with generators \(x^p\) and \(y^q\),
\(\log|x^p|(t)=-pt\) and \(\log|y^q|(t)=-q(1-t)\), and the maximum of the
two logarithms is taken.  Second, the breakpoint is where the two branches
of the piecewise affine function cross, and the divisor \(D\) with
\(\mathfrak a\cO_{\mathcal Y}=\cO_{\mathcal Y}(D)\) lives on the model
\(\mathcal Y\) obtained by subdividing the edge there: for
\(p=q=1\) the breakpoint is \(t=1/2\), which is not a lattice point of the
original edge, so \(\mathcal Y\) is a genuine blowup and
\(\max\{\varphi_1,\varphi_2\}\) is \emph{not} determined on the original
model.  This is the one-dimensional shadow of
\Cref{ex:ch20-max-fails-naively}, and it is why
\Cref{prop:ch20-curve-lattice}, which works with a fixed model, speaks of
vertex vectors and not of arbitrary maxima.
\end{example}

\begin{example}[Toric support functions]
\label{ex:ch20-toric}
On a toric model, a vertical Cartier divisor determines a rational
piecewise affine support function.  Passing to a toric subdivision is
pullback, adding a vertical divisor adds its support function, and relative
nefness is convexity after the reference support function has been added.
Effective divisors take nonnegative values at their divisorial points by
\eqref{eq:ch20-sign}.  This simultaneously checks the signs of
\(\ddc\varphi_D=[D]\) and of the maximum construction.

Made explicit, the dictionary is the following.  Let \(X\) be a polarized
toric \(K\)-variety with polytope \(P\subseteq M_{\mathbf R}\), and use the
tropical coordinate \(u=-\log|\chi|\) of \eqref{eq:appJ-toric-coordinate}.
An invariant closed model form \(\theta\) with the reference model metric
corresponds to the recession datum \(\check h_P\) of
\eqref{eq:appJ-toric-recession}, an invariant model function corresponds to
a rational piecewise affine \(\psi\) on \(N_{\mathbf R}\) with bounded
difference from \(0\), and
\begin{itemize}
\item passing to a toric refinement of the fan is pullback along a
 domination, and does not change \(\psi\);
\item adding a vertical divisor adds its piecewise affine function, matching
 \(\ddc\psi\) with the passage from \(\psi\) to its associated class;
\item \(\theta+\ddc\psi\) is semipositive if and only if
 \(\check h_P+\psi\) is convex, which is the toric form of
 \Cref{def:ch20-semipositive}.
\end{itemize}
In this language \Cref{prop:ch20-model-max} becomes the elementary statement
that a maximum of two convex functions is convex, and
\Cref{ex:ch20-max-fails-naively} becomes the equally elementary statement
that a maximum of two functions affine on the cells of a given fan is
piecewise affine for a \emph{refinement} of that fan---the refinement being
the normalized blowup of \eqref{eq:ch20-ideal-divisor}.  The same
description reappears in \Cref{chap:bfj-psh} for invariant psh functions.

The one number worth carrying away is the total mass.  With the lattice
normalization \eqref{eq:appJ-alexandrov} and the factor \(n!\) of
\eqref{eq:appJ-toric-degree}, the Monge--Amp\`ere mass of an invariant
semipositive model metric on \((X,L)\) is
\(n!\operatorname{vol}_M(P)=(L^n)\), in agreement with
\eqref{eq:ch20-special-fibre-degree} and with
\Cref{thm:toric-ma-comparison}; for \(X=\mathbf P^1\), \(L=\cO(d)\) and
\(P=[0,d]\), this reads \(1!\cdot d=d\).  Effective divisors take
nonnegative values at their divisorial points by \eqref{eq:ch20-sign}, which
fixes the remaining sign.  A reader comparing with a source that uses
\(\log|\chi|\) rather than \(-\log|\chi|\) as tropical coordinate should
first apply \eqref{eq:appJ-toric-recession}; the subgradient image lies in
\(-P\), not \(P\).
\end{example}

\section{Scope: the original BFJ hypotheses and later extensions}
\label{sec:ch20-scope}

\Cref{part:bfj} is written under the signature
\eqref{eq:ch20-standing}, and this section records, statement by statement,
which parts of that signature are actually used and what is known without
them.  The purpose is not completeness but auditability: a reader who
wishes to apply a result of this chapter over a non-discretely valued field,
in positive residue characteristic, or on a singular \(X\) should be able to
see at once whether the proof given here survives.

\subsection{Where each hypothesis enters}
\label{sec:ch20-hypothesis-use}

\begin{center}
\footnotesize
\begin{tabular}{p{.26\textwidth}|p{.20\textwidth}|p{.46\textwidth}}
\toprule
\textbf{Statement}&\textbf{Uses}&\textbf{Comment}\\
\midrule
\Cref{lem:ch20-special-fibre-test}
 &\(R\) local; \(\cX\) projective over \(R\)
 &Discreteness is not used; any valuation ring with local spectrum works,
  provided \cite[4.7.1]{EGAIII1} applies, that is, provided \(R\) is
  noetherian.  For a non-noetherian valuation ring one works with
  admissible formal models instead \textup{(\Cref{chap:models})}.\\
\Cref{lem:ch20-pullback-injective},
\Cref{prop:ch20-two-model-comparison}
 &directedness of models
 &Formal; uses only \Cref{prop:common-domination} and the projection
  formula.  Valid for any complete nontrivially valued \(K\) and any normal
  projective \(X\).\\
\Cref{prop:ch20-semipositive-well-defined}
 &projection formula
 &Valid over any complete nontrivially valued field; this is the
  model-metric statement \Cref{lem:metrics-nef-independence}.\\
\Cref{lem:ch20-ample-propagation}
 &\(\cX,\cY\) projective over \(R\)
 &Uses only \Cref{thm:appB-admissible-domination} and
  \cite[4.6.13]{EGAII}; no regularity, no residue-characteristic
  hypothesis.\\
\Cref{prop:ch20-positive-representative}
 &\(X\) projective; \(\{\theta\}\) ample; plus
  \Cref{thm:ch20-ddc-exact}
 &Step~1 uses the flat-closure construction of \Cref{chap:models}, written
  for \(K\) discretely valued; Step~2 uses exactness of
  \eqref{eq:ch20-ddc-exact} in the direct limit, and therefore inherits its
  hypotheses.  Smoothness of \(X\) is used only through that.\\
\Cref{thm:ch20-ddc-exact}
 &\(K\) discretely valued; \(\operatorname{char}k=0\); \(X\) smooth and
  connected
 &Regular models must be cofinal, which is resolution of singularities;
  connectedness and normality give the connected special fibre; regularity
  makes each \(E_i\) Cartier and \(\operatorname{CH}^1=\Pic\).  This is the
  most hypothesis-sensitive statement of the chapter.\\
\Cref{prop:ch20-model-max}
 &\(\cX\) projective over \(R\); \(N^1(\cX/R)\) spanned by ample classes
 &Neither regularity nor \(\operatorname{char}k=0\) is used.  The
  corresponding statement for model metrics over an arbitrary complete
  field is \Cref{prop:semipositive-stability}\,(3).\\
\Cref{thm:ch20-vertical-nef-rigidity}
 &regular \(\cX\); connected \(\cX_s\)
 &Uses the equality case of \eqref{eq:ch20-hodge-laplacian}, hence
  regularity, but not \Cref{lem:ch20-kunnemann-input}.\\
\Cref{prop:ch20-relative-hodge},
\Cref{thm:ch20-mixed-hodge}
 &regular \(\cX\); \(\cX_s\) connected
 &Regularity is used only so that the \(E_i\) are Cartier and the
  intersection form is symmetric; for a normal \(\cX\) one restricts to
  \(\mathbf R\)-Cartier vertical divisors.  Residue characteristic is
  irrelevant.\\
\Cref{prop:ch20-curve-lattice},
\Cref{prop:ch20-curve-envelope}
 &\(n=1\); regular \(\cX\)
 &Genuinely one-dimensional; see
  \Cref{rem:ch20-fixed-model-envelope-scope}.\\
\bottomrule
\end{tabular}
\end{center}

\subsection{What later work removes, and what it does not}
\label{sec:ch20-later-extensions}

Three families of extensions are relevant to this chapter, and each is
labelled as such wherever it is used later in the book.

\emph{(E1) Arbitrary complete nontrivially valued fields.}  Gubler and
Martin \cite{GM19} compare algebraic, formal, and piecewise linear metrics
and their semipositivity over an arbitrary complete nontrivially valued
field, without discreteness and without any hypothesis on the residue
characteristic; their comparison is quoted in this book as
\Cref{thm:gm-metric-comparison}.  What transfers immediately is the
definition of a model metric, the well-definedness of semipositivity
(\Cref{lem:metrics-nef-independence}), and the maximum construction.  What
does \emph{not} transfer without further work is
\Cref{thm:ch20-ddc-exact}: over a non-discretely valued field the special
fibre of an admissible formal model is not a Cartier divisor with a
canonical multiplicity vector, so both the statement
\(\varphi_{\cX_s}=1\) and the row-sum relation
\(\sum_jb_jc_{ij}=0\) have to be reformulated before the argument can be
attempted.  We do not attempt it.

\emph{(E2) Positive residue characteristic.}  Resolution of singularities is
unavailable, so regular models need not be cofinal and
\Cref{prop:snc-cofinal} fails as stated; only alterations are available
(\Cref{app:birational-geometry}).  All statements of this chapter that do
not mention regularity survive verbatim, in particular
\Cref{lem:ch20-special-fibre-test,lem:ch20-pullback-injective,prop:ch20-two-model-comparison,prop:ch20-semipositive-well-defined,prop:ch20-model-max}.
The results that do not survive as proved here are
\Cref{thm:ch20-ddc-exact} and its consequences
\Cref{cor:ch20-dimension-count}, \Cref{cor:ch20-normal-form} and
\Cref{prop:ch20-positive-representative}.  On any \emph{fixed regular}
model the whole of Step~1 of \Cref{thm:ch20-ddc-exact} remains valid, and
with it \Cref{thm:ch20-vertical-nef-rigidity}, the inequalities
\eqref{eq:ch20-relative-hodge}, \eqref{eq:ch20-mixed-hodge},
\eqref{eq:ch20-mixed-cauchy-schwarz} and their equality clauses; what is
lost is the passage to the direct limit, because regular models need not be
cofinal.  The Künnemann step, Step~4, is insensitive to the residue
characteristic.  In
\Cref{chap:bfj-compactness,chap:bfj-envelopes} the role of residue
characteristic zero is different and much deeper: it is needed for
multiplier ideals, and the substitute in positive characteristic is the
test-ideal technique of Gubler, Jell, K\"unnemann and Martin
\cite{GJKM19}.  Nothing in the present chapter depends on that.

\emph{(E3) Singular and non-smooth \(X\).}  Smoothness of \(X\) is used in
this chapter only through connectedness and normality of the models, in
Step~1 of \Cref{thm:ch20-ddc-exact}.  A normal projective \(X\) suffices
for every statement proved here.  The later chapters use smoothness more
seriously---for instance, through the existence of SNC models with a
prescribed skeleton---and the exact place where it becomes essential is
recorded in \Cref{chap:bfj-compactness}.

\begin{warning}{The scope table is about this chapter's proofs}
The entries above describe what the arguments \emph{given here} require.
They are not claims that the corresponding statements are false under weaker
hypotheses, nor that they are unknown.  Where a published extension exists
it is named; where none is named, the honest reading is ``not proved in this
book''.  This is the discipline announced in the introduction to
\Cref{part:bfj}, and it is applied uniformly in
\Crefrange{chap:bfj-psh}{chap:bfj-examples}.
\end{warning}

\section{Dictionary}
\label{sec:ch20-dictionary}

\subsection{Objects and their two descriptions}
\label{sec:ch20-dict-objects}

\begin{center}
\small
\begin{tabular}{p{.29\textwidth}|p{.30\textwidth}|p{.33\textwidth}}
\toprule
\textbf{Algebraic datum}&\textbf{Pluripotential object}
 &\textbf{Established at}\\
\midrule
relative numerical class on \(\cX\)
 &closed model \((1,1)\)-form determined on \(\cX\)
 &\Cref{def:ch20-closed-model-form}\\
class pulled back to a common domination
 &equality of two closed model forms
 &\Cref{prop:ch20-two-model-comparison}\\
vertical \(\mathbf R\)-Cartier divisor \(D\)
 &model function \(\varphi_D=v(f_D)\)
 &\eqref{eq:ch20-sign},
  \Cref{prop:metrics-divisor-model-function}\\
numerical class \([D]\) of a vertical divisor
 &\(\ddc\varphi_D\)
 &\Cref{def:ch20-ddc},
  \Cref{lem:ch20-ddc-well-defined}\\
multiple of the special fibre
 &additive constant
 &\eqref{eq:ch20-ddc-constant},
  \Cref{thm:ch20-ddc-exact}\\
relative nefness of \(\theta_{\cX}\)
 &semipositivity of \(\theta\)
 &\Cref{def:ch20-semipositive},
  \Cref{prop:ch20-semipositive-well-defined}\\
relative ampleness of \(\theta_{\cX}\)
 &\(\cX\)-positivity of \(\theta\)
 &\Cref{def:ch20-four-positivities},
  \Cref{ex:ch20-nef-not-ample}\\
sum of vertical fractional ideals
 &maximum of model functions
 &\eqref{eq:ch20-ideal-sum-max},
  \Cref{prop:ch20-model-max}\\
degree of \(D\) on the components
 &weighted graph Laplacian of \(\varphi_D\)
 &\Cref{prop:ch20-curve-laplacian}\\
vertical intersection number
 &integral against a Chambert--Loir measure
 &\eqref{eq:ch20-pairing-as-integral},
  \Cref{prop:ch20-measure-positivity}\\
mixed vertical sum
 \(\frac1{n+1}\sum_{j=0}^n(D\cdot\theta_1^j\cdot\theta_0^{n-j})\),
 where \(\theta_1=\theta_0+\ddc\varphi_D\)
 &energy difference between cohomologous model forms
 &\Cref{prop:ch20-intersection-pairing},
  \Cref{cor:ch20-concavity}\\
negativity of \((D^2\cdot\alpha_1\cdots\alpha_{n-1})\)
 &concavity of energy; Cauchy--Schwarz
 &\Cref{thm:ch20-mixed-hodge}\\
\bottomrule
\end{tabular}
\end{center}

\subsection{Reading the dictionary backwards}
\label{sec:ch20-dict-backward}

Four entries above are one-way, and it is worth saying which.

\emph{From a model function to a divisor} is well defined only after a model
has been chosen: the divisor is unique on a fixed normal model
(\Cref{prop:divisor-model-function-injective-mod-fibre}) but changes under
pullback.  \emph{From a closed model form to a determination} is likewise
model dependent, and the only invariant statements are those proved to be
independent of the determination.  \emph{From semipositivity to relative
ampleness} there is no arrow at all, only the repair
\Cref{prop:ch20-positive-representative}, which changes the form by
\(\ddc\psi\).  And \emph{from a nef generic class to a semipositive form}
there is no arrow either, by
\Cref{thm:ch20-vertical-nef-rigidity}: the vertical directions carry no
positivity of their own.

\subsection{Numbers to remember}
\label{sec:ch20-dict-numbers}

The following values were computed in full and are the test suite for the
rest of \Cref{part:bfj}.

\begin{center}
\small
\begin{tabular}{p{.34\textwidth}|p{.26\textwidth}|p{.32\textwidth}}
\toprule
\textbf{Quantity}&\textbf{Value}&\textbf{Where}\\
\midrule
degrees of \(\theta+\ddc\varphi\) on the running conic
 &\((1+B-A,\ 1+A-B)\)
 &\eqref{eq:ch20-conic-perturbed-degrees}\\
semipositivity range on the conic
 &\(|A-B|\leq1\)
 &\Cref{ex:ch20-conic-ddc},
  \eqref{eq:running-nef-interval}\\
\(\dim N^1(\cX/R)\) for \(\mathbf P^2_R\) and its blowups
 &\(1,\ 2,\ 3\)
 &\eqref{eq:ch20-two-model-dimensions}\\
\((E^3)\) on \(\operatorname{Bl}_p\mathbf P^2_R\)
 &\(+1\)
 &\eqref{eq:ch20-triple-basic}\\
\(3E_{\theta_0}(t\varphi)\) for
 \(\theta_t=2\rho^*h-t[E]\), \(\varphi=-\varphi_E\)
 &\(-t^3\), concave on \([0,2]\)
 &\eqref{eq:ch20-triple-energy}\\
\((D^2\cdot\mathcal A)\) for \(D=aE\)
 &\(-a^2\)
 &\eqref{eq:ch20-triple-hodge-check}\\
breakpoint of \(\log|(x^p,y^q)|\)
 &\(t=q/(p+q)\)
 &\eqref{eq:ch20-nodal-ideal}\\
\bottomrule
\end{tabular}
\end{center}

\section*{Notes and further reading}

The construction of model functions, closed forms, the exact sequence
\eqref{eq:ch20-ddc-exact}, positivity, and the maximum theorem follows
\cite[Sections~2, 4, and~5]{BFJ16}.  The exactness proof uses the relative
Hodge argument and K\"unnemann's extension result; the latter is a genuine
algebraic input, not a consequence of the direct-limit notation.  Model
metrics over more general complete fields and the comparison among
algebraic, formal, and piecewise linear semipositivity are treated in
\cite{GM19}.  The sign convention is always
\eqref{eq:ch20-sign}; no complex normalization of \(\ddc\) is being used.

\paragraph{Provenance, statement by statement.}
The direct-limit definition of \(\mathcal Z^{1,1}(X)\), the model
differential, and the exact sequence are \cite[\S\S2 and 4]{BFJ16}, the
source referred to throughout \Cref{part:bfj}; the correspondence input
isolated here as \Cref{lem:ch20-kunnemann-input} is K\"unnemann's, and the
relative Hodge index argument in the form
\eqref{eq:ch20-hodge-laplacian} goes back to the classical theory of
arithmetic surfaces.  The maximum theorem
\Cref{prop:ch20-model-max} and the ideal-theoretic device
\eqref{eq:ch20-max-ideal} are \cite[\S5]{BFJ16}, and
\Cref{prop:ch20-positive-representative} is the model-level positivity
statement of \cite[\S1]{BFJ16}, here given a proof that separates the
propagation of relative ampleness (\Cref{lem:ch20-ample-propagation}) from
the exactness input.  The non-Archimedean
Monge--Amp\`ere equation solved by the same authors in \cite{BFJ15} is the
principal application, and its survey \cite{BFJ16b} is a short entry point
to the circle of ideas.  Zhang's notion of a semipositive metric
\cite{Zhang1995} is the ancestor of \Cref{def:ch20-semipositive}, and
\Cref{lem:metrics-nef-independence} is the statement that the two agree on
model metrics.

\paragraph{What is organized differently here.}
Three choices depart from the arrangement in the sources.  First, we have
separated the \emph{construction} of the direct limit
(\Cref{sec:ch20-numerical,sec:ch20-model-forms}) from the \emph{exactness
theorem} (\Cref{sec:ch20-exact}), so that the two-model comparison
criterion \Cref{prop:ch20-two-model-comparison}---which is what a reader
actually uses when computing---is available before any deep input.  Second,
the four positivity notions have been made into a single definition with a
single theorem recording all implications and all failures
(\Cref{def:ch20-four-positivities},
\Cref{thm:ch20-positivity-implications},
\Cref{thm:ch20-vertical-nef-rigidity}); in the literature these are
scattered, and the rigidity statement, though elementary, is rarely made
explicit.  Third, all examples are carried to numbers and cross-checked
against \Cref{chap:metrics,chap:model-calculations}: the running conic in
\Cref{ex:ch20-conic-ddc}, the blown-up plane in
\Cref{ex:ch20-p2-blowup-numerics,ex:ch20-triple-products}, and the two-model
comparison in \Cref{ex:ch20-two-model-surface}.  The last has no counterpart
in the sources and is the smallest example in which the direct limit does
something visible.

\paragraph{Interfaces.}
The comparison with the Chambert--Loir--Ducros calculus is
\Cref{thm:ch19-cld-bfj-comparison}, proved in
\Cref{chap:cld-examples}; its scope clause, and the limitations table of
\Cref{sec:ch19-limitations}, delimit exactly what may and may not be
transported between the two theories.  The Chambert--Loir measures used in
\eqref{eq:ch20-pairing-as-integral} were constructed in
\Cref{sec:metrics-cl-measure}, and their identification with a wedge product
of curvature currents is
\Cref{thm:cld-recovers-cl}.  The trivially valued analogue of the present
chapter is \Cref{chap:bj-fs-pl-psh}, where the inverse system of models is
replaced by the system of birational models of \(X\) over the ground field
and the special fibre by the exceptional locus; the resulting formalism is
similar but is not obtained from this one by setting the norm equal to one,
as explained in the introduction to \Cref{part:bj}.

\paragraph{What is new here, and what is not.}
No theorem of this chapter is new.  \Cref{thm:ch20-vertical-nef-rigidity},
\Cref{prop:ch20-curve-lattice}, \Cref{prop:ch20-curve-envelope} and the
numerical verifications of \Cref{sec:ch20-examples} are elementary
consequences of statements in the sources; we have written them out because
each of them is used later in a form that is easy to misremember, and
because a chapter whose entire content is a dictionary is worth only as much
as the smallest example in which both sides have been computed.
